
\documentclass{article}
\usepackage{colortbl}
\usepackage{makecell}
\usepackage{multirow}
\usepackage{supertabular}

\begin{document}

\newcounter{utterance}

\centering \large Interaction Transcript for game `clean\_up', experiment `1\_medium\_5obj\_ar', episode 2 with gemma{-}4{-}26b{-}a4b{-}it.
\vspace{24pt}

{ \footnotesize  \setcounter{utterance}{1}
\setlength{\tabcolsep}{0pt}
\begin{supertabular}{c@{$\;$}|p{.15\linewidth}@{}p{.15\linewidth}p{.15\linewidth}p{.15\linewidth}p{.15\linewidth}p{.15\linewidth}}
    \# & $\;$A & \multicolumn{4}{c}{Game Master} & $\;\:$B\\
    \hline

    \theutterance \stepcounter{utterance}  
    & & \multicolumn{4}{p{0.6\linewidth}}{
        \cellcolor[rgb]{0.9,0.9,0.9}{
            \makecell[{{p{\linewidth}}}]{
                \texttt{\tiny{[P1$\langle$GM]}}
                \texttt{أنا مدير لعبتك، وأنت تلعب لعبة تعاونية مع الشبكة التالية كلوحة لعب:} \\
\\ 
\texttt{```} \\
\texttt{ 1234567} \\
\texttt{╔═╤═╤═══╗} \\
\texttt{║◌│◌│◌◌◌║ 1} \\
\texttt{╟─┴─┼─┬─╢ 2} \\
\texttt{║◌◌◌│◌│T║ 3} \\
\texttt{║◌H◌├─┘◌║ 4} \\
\texttt{║◌┌─┘◌◌W║ 5} \\
\texttt{╟─┤◌◌◌◌◌║ 6} \\
\texttt{║I│◌C◌◌◌║ 7} \\
\texttt{╚═╧═════╝} \\
\\ 
\texttt{```} \\
\\ 
\texttt{* تعرض الحافة العلوية إحداثيات x التي تتزايد نحو اليمين، والحافة اليمنى إحداثيات y التي تتزايد نحو الأسفل.} \\
\texttt{* الكائنات التالية موضوعة عشوائياً على شبكتك: 'W', 'I', 'T', 'C', 'H'.} \\
\\ 
\\ 
\texttt{يرى اللاعب الآخر نسخة مختلفة من لوحة اللعب، حيث توضع الكائنات في مواقع عشوائية مختلفة. لا يمكنك رؤية لوحة اللاعب الآخر، ولا يمكنه رؤية لوحتك.} \\
\\ 
\\ 
\texttt{**الهدف:**} \\
\\ 
\texttt{يجب على كلا اللاعبين تحريك الكائنات على خلفية كل منهما بحيث تنتهي الكائنات المتطابقة عند الإحداثيات نفسها. عليك التواصل مع اللاعب الآخر للاتفاق على إعداد هدف مشترك.} \\
\\ 
\\ 
\texttt{**القواعد:**} \\
\\ 
\texttt{* في كل دور، يمكنك إرسال أمر واحد فقط من الأمرين التاليين:} \\
\texttt{ 1. `قل: <رسالة>`: لإرسال رسالة (كل شيء حتى فاصل السطر التالي) إلى اللاعب الآخر. سأقوم بتمريرها إلى شريكك.} \\
\texttt{2. `حرّك: <كائن>, (<X>, <Y>)`: لتحريك كائن إلى موضع جديد، حيث `<X>` هو العمود و `<Y>` هو الصف. سأخبرك بما إذا كانت حركتك ناجحة أم لا.} \\
\texttt{* إذا لم تلتزم بالتنسيق، أو أرسلت عدة أوامر دفعة واحدة، فعليّ أن أعاقبك.} \\
\texttt{* إذا تجاوز كلا اللاعبين أكثر من 12 عقوبة، فستخسران اللعبة معاً.} \\
\texttt{* من الضروري أن تتواصل مع اللاعب الآخر بشأن حالة هدفك! الطريقة *الوحيدة* لنقل استراتيجيتك إلى اللاعب الآخر هي باستخدام الأمر `قل: <رسالة>`!} \\
\\ 
\\ 
\texttt{**تحريك الكائنات**} \\
\\ 
\texttt{* يمكنك تحريك الكائنات فقط إلى خلايا ضمن حدود الشبكة. يجب أن تكون الخلية المستهدفة فارغة، أي يجب أن تحتوي فقط على الرمز '◌'.} \\
\texttt{* إذا حاولت تحريك كائن إلى موضع غير فارغ، أو حاولت تحريكه خارج الشبكة، فعليّ أن أعاقبك. ستحصل على محاولة أخرى.} \\
\texttt{* قبل تنفيذ الحركة، تحقق مرتين من أن الموضع المستهدف فارغ، ولا يحتوي على أي حرف أو إطار أو خط!} \\
\\ 
\\ 
\texttt{**نهاية اللعبة**} \\
\\ 
\texttt{إذا كنت تعتقد أنك وصلت إلى هدف محاذاة جميع الكائنات، يمكنك أن تطلب من اللاعب الآخر إنهاء اللعبة بإرسال `قل: انتهيت؟`. إذا طلب منك اللاعب الآخر إنهاء اللعبة، ورددت بـ `قل: انتهيت!`، ستنتهي اللعبة.} \\
\\ 
\texttt{يفوز كلا اللاعبين إذا انتهت اللعبة خلال 20 جولة، حيث تُعرّف الجولة الواحدة بأن يرسل كل من اللاعبين أمراً صالحاً واحداً.} \\
\\ 
\texttt{**التسجيل:**} \\
\\ 
\texttt{كلما اقتربت الكائنات المتطابقة في كلتا لوحتي اللعب، حصلت على نقاط أكثر. العقوبات تقلل نقاطك. هل يمكنك تحطيم الرقم القياسي؟} \\
\\ 
\\ 
\texttt{يرجى إرسال رسالة إلى اللاعب الآخر لبدء اللعبة!} \\
            }
        }
    }
    & \\ \\

    \theutterance \stepcounter{utterance}  
    & \multicolumn{4}{p{0.6\linewidth}}{
        \cellcolor[rgb]{0.9,0.9,0.9}{
            \makecell[{{p{\linewidth}}}]{
                \texttt{\tiny{[P1$\rangle$GM]}}
                \texttt{قل: مرحباً! لنبدأ بالتنسيق. سأقوم بإرسال مواقع الكائنات الموجودة لدي الآن لنحاول إيجاد هدف مشترك.} \\
\\ 
\texttt{مواقع الكائنات لدي هي:} \\
\texttt{W: (6, 5)} \\
\texttt{I: (1, 7)} \\
\texttt{T: (6, 3)} \\
\texttt{C: (4, 7)} \\
\texttt{H: (2, 4)} \\
\\ 
\texttt{من فضلك، أرسل لي مواقع الكائنات لديك لنحدد الإحداثيات التي سنحاول الوصول إليها جميعاً.} \\
            }
        }
    }
    & & \\ \\

    \theutterance \stepcounter{utterance}  
    & & & \multicolumn{2}{p{0.3\linewidth}}{
        \cellcolor[rgb]{0.9,0.9,0.9}{
            \makecell[{{p{\linewidth}}}]{
                \texttt{\tiny{[GM$|$GM]}}
                \texttt{قل: مرحباً! لنبدأ بالتنسيق. سأقوم بإرسال مواقع الكائنات الموجودة لدي الآن لنحاول إيجاد هدف مشترك.} \\
\\ 
\texttt{مواقع الكائنات لدي هي:} \\
\texttt{W: (6, 5)} \\
\texttt{I: (1, 7)} \\
\texttt{T: (6, 3)} \\
\texttt{C: (4, 7)} \\
\texttt{H: (2, 4)} \\
\\ 
\texttt{من فضلك، أرسل لي مواقع الكائنات لديك لنحدد الإحداثيات التي سنحاول الوصول إليها جميعاً.} \\
            }
        }
    }
    & & \\ \\

    \theutterance \stepcounter{utterance}  
    & & & \multicolumn{2}{p{0.3\linewidth}}{
        \cellcolor[rgb]{0.9,0.9,0.9}{
            \makecell[{{p{\linewidth}}}]{
                \texttt{\tiny{[GM$|$GM]}}
                \texttt{Invalid format: tail is not empty: '} \\
\\ 
\texttt{مواقع الكائنات لدي هي:} \\
\texttt{W: (6, 5)} \\
\texttt{I: (1, 7)} \\
\texttt{T: (6, 3)} \\
\texttt{C: (4, 7)} \\
\texttt{H: (2, 4)} \\
\\ 
\texttt{من فضلك، أرسل لي مواقع الكائنات لديك لنحدد الإحداثيات التي سنحاول الوصول إليها جميعاً.} \\
            }
        }
    }
    & & \\ \\

    \theutterance \stepcounter{utterance}  
    & & \multicolumn{4}{p{0.6\linewidth}}{
        \cellcolor[rgb]{0.9,0.9,0.9}{
            \makecell[{{p{\linewidth}}}]{
                \texttt{\tiny{[P1$\langle$GM]}}
                \texttt{عقوبة: يجب ألا تحتوي رسالتك على أي شيء بعد الأمر!} \\
\\ 
\texttt{تأكد من أن ردك يحتوي فقط على `قل: <رسالة>` أو `حرّك: <كائن>, (<X>, <Y>)`، ولا شيء غير ذلك!} \\
\\ 
\texttt{لقد جمعتما معاً 1 من أصل 12 عقوبة.} \\
\texttt{يرجى المحاولة مرة أخرى!} \\
            }
        }
    }
    & \\ \\

    \theutterance \stepcounter{utterance}  
    & \multicolumn{4}{p{0.6\linewidth}}{
        \cellcolor[rgb]{0.9,0.9,0.9}{
            \makecell[{{p{\linewidth}}}]{
                \texttt{\tiny{[P1$\rangle$GM]}}
                \texttt{قل: مرحباً! لنبدأ بالتنسيق. سأقوم بإرسال مواقع الكائنات الموجودة لدي الآن لنحاول إيجاد هدف مشترك. مواقع الكائنات لدي هي: W: (6, 5), I: (1, 7), T: (6, 3), C: (4, 7), H: (2, 4). من فضلك، أرسل لي مواقع الكائنات لديك لنحدد الإحداثيات التي سنحاول الوصول إليها جميعاً.} \\
            }
        }
    }
    & & \\ \\

    \theutterance \stepcounter{utterance}  
    & & & \multicolumn{2}{p{0.3\linewidth}}{
        \cellcolor[rgb]{0.9,0.9,0.9}{
            \makecell[{{p{\linewidth}}}]{
                \texttt{\tiny{[GM$|$GM]}}
                \texttt{قل: مرحباً! لنبدأ بالتنسيق. سأقوم بإرسال مواقع الكائنات الموجودة لدي الآن لنحاول إيجاد هدف مشترك. مواقع الكائنات لدي هي: W: (6, 5), I: (1, 7), T: (6, 3), C: (4, 7), H: (2, 4). من فضلك، أرسل لي مواقع الكائنات لديك لنحدد الإحداثيات التي سنحاول الوصول إليها جميعاً.} \\
            }
        }
    }
    & & \\ \\

    \theutterance \stepcounter{utterance}  
    & & \multicolumn{4}{p{0.6\linewidth}}{
        \cellcolor[rgb]{0.9,0.9,0.9}{
            \makecell[{{p{\linewidth}}}]{
                \texttt{\tiny{[GM$\rangle$P2]}}
                \texttt{أنا مدير لعبتك، وأنت تلعب لعبة تعاونية مع الشبكة التالية كلوحة لعب:} \\
\\ 
\texttt{```} \\
\texttt{ 1234567} \\
\texttt{╔═╤═╤═══╗} \\
\texttt{║◌│◌│◌◌◌║ 1} \\
\texttt{╟─┴─┼─┬─╢ 2} \\
\texttt{║◌◌◌│◌│◌║ 3} \\
\texttt{║◌◌C├─┘◌║ 4} \\
\texttt{║◌┌─┘◌T◌║ 5} \\
\texttt{╟─┤◌W◌◌H║ 6} \\
\texttt{║◌│◌I◌◌◌║ 7} \\
\texttt{╚═╧═════╝} \\
\\ 
\texttt{```} \\
\\ 
\texttt{* تعرض الحافة العلوية إحداثيات x التي تتزايد نحو اليمين، والحافة اليمنى إحداثيات y التي تتزايد نحو الأسفل.} \\
\texttt{* الكائنات التالية موضوعة عشوائياً على شبكتك: 'W', 'I', 'T', 'C', 'H'.} \\
\\ 
\\ 
\texttt{يرى اللاعب الآخر نسخة مختلفة من لوحة اللعب، حيث توضع الكائنات في مواقع عشوائية مختلفة. لا يمكنك رؤية لوحة اللاعب الآخر، ولا يمكنه رؤية لوحتك.} \\
\\ 
\\ 
\texttt{**الهدف:**} \\
\\ 
\texttt{يجب على كلا اللاعبين تحريك الكائنات على خلفية كل منهما بحيث تنتهي الكائنات المتطابقة عند الإحداثيات نفسها. عليك التواصل مع اللاعب الآخر للاتفاق على إعداد هدف مشترك.} \\
\\ 
\\ 
\texttt{**القواعد:**} \\
\\ 
\texttt{* في كل دور، يمكنك إرسال أمر واحد فقط من الأمرين التاليين:} \\
\texttt{ 1. `قل: <رسالة>`: لإرسال رسالة (كل شيء حتى فاصل السطر التالي) إلى اللاعب الآخر. سأقوم بتمريرها إلى شريكك.} \\
\texttt{2. `حرّك: <كائن>, (<X>, <Y>)`: لتحريك كائن إلى موضع جديد، حيث `<X>` هو العمود و `<Y>` هو الصف. سأخبرك بما إذا كانت حركتك ناجحة أم لا.} \\
\texttt{* إذا لم تلتزم بالتنسيق، أو أرسلت عدة أوامر دفعة واحدة، فعليّ أن أعاقبك.} \\
\texttt{* إذا تجاوز كلا اللاعبين أكثر من 12 عقوبة، فستخسران اللعبة معاً.} \\
\texttt{* من الضروري أن تتواصل مع اللاعب الآخر بشأن حالة هدفك! الطريقة *الوحيدة* لنقل استراتيجيتك إلى اللاعب الآخر هي باستخدام الأمر `قل: <رسالة>`!} \\
\\ 
\\ 
\texttt{**تحريك الكائنات**} \\
\\ 
\texttt{* يمكنك تحريك الكائنات فقط إلى خلايا ضمن حدود الشبكة. يجب أن تكون الخلية المستهدفة فارغة، أي يجب أن تحتوي فقط على الرمز '◌'.} \\
\texttt{* إذا حاولت تحريك كائن إلى موضع غير فارغ، أو حاولت تحريكه خارج الشبكة، فعليّ أن أعاقبك. ستحصل على محاولة أخرى.} \\
\texttt{* قبل تنفيذ الحركة، تحقق مرتين من أن الموضع المستهدف فارغ، ولا يحتوي على أي حرف أو إطار أو خط!} \\
\\ 
\\ 
\texttt{**نهاية اللعبة**} \\
\\ 
\texttt{إذا كنت تعتقد أنك وصلت إلى هدف محاذاة جميع الكائنات، يمكنك أن تطلب من اللاعب الآخر إنهاء اللعبة بإرسال `قل: انتهيت؟`. إذا طلب منك اللاعب الآخر إنهاء اللعبة، ورددت بـ `قل: انتهيت!`، ستنتهي اللعبة.} \\
\\ 
\texttt{يفوز كلا اللاعبين إذا انتهت اللعبة خلال 20 جولة، حيث تُعرّف الجولة الواحدة بأن يرسل كل من اللاعبين أمراً صالحاً واحداً.} \\
\\ 
\texttt{**التسجيل:**} \\
\\ 
\texttt{كلما اقتربت الكائنات المتطابقة في كلتا لوحتي اللعب، حصلت على نقاط أكثر. العقوبات تقلل نقاطك. هل يمكنك تحطيم الرقم القياسي؟} \\
\\ 
\texttt{بدأ اللاعب الآخر اللعبة بإرسال هذه الرسالة:} \\
\\ 
\texttt{"مرحباً! لنبدأ بالتنسيق. سأقوم بإرسال مواقع الكائنات الموجودة لدي الآن لنحاول إيجاد هدف مشترك. مواقع الكائنات لدي هي: W: (6, 5), I: (1, 7), T: (6, 3), C: (4, 7), H: (2, 4). من فضلك، أرسل لي مواقع الكائنات لديك لنحدد الإحداثيات التي سنحاول الوصول إليها جميعاً."} \\
\\ 
\texttt{ما هو أمرك الأول؟} \\
            }
        }
    }
    & \\ \\

    \theutterance \stepcounter{utterance}  
    & & & \multicolumn{4}{p{0.6\linewidth}}{
        \cellcolor[rgb]{0.9,0.9,0.9}{
            \makecell[{{p{\linewidth}}}]{
                \texttt{\tiny{[GM$\langle$P2]}}
                \texttt{قل: مرحباً! لقد استلمت مواقع الكائنات لديك. إليك مواقع الكائنات الموجودة في لوحتي حالياً: W: (3, 6), I: (3, 7), T: (6, 5), C: (3, 4), H: (7, 6). لنحاول الاتفاق على إحداثيات هدف موحدة.} \\
            }
        }
    }
     \\ \\

    \theutterance \stepcounter{utterance}  
    & & & \multicolumn{2}{p{0.3\linewidth}}{
        \cellcolor[rgb]{0.9,0.9,0.9}{
            \makecell[{{p{\linewidth}}}]{
                \texttt{\tiny{[GM$|$GM]}}
                \texttt{قل: مرحباً! لقد استلمت مواقع الكائنات لديك. إليك مواقع الكائنات الموجودة في لوحتي حالياً: W: (3, 6), I: (3, 7), T: (6, 5), C: (3, 4), H: (7, 6). لنحاول الاتفاق على إحداثيات هدف موحدة.} \\
            }
        }
    }
    & & \\ \\

    \theutterance \stepcounter{utterance}  
    & & \multicolumn{4}{p{0.6\linewidth}}{
        \cellcolor[rgb]{0.9,0.9,0.9}{
            \makecell[{{p{\linewidth}}}]{
                \texttt{\tiny{[P1$\langle$GM]}}
                \texttt{تم تمرير رسالتك إلى اللاعب الآخر.} \\
\\ 
\\ 
\texttt{أنت تلعب حالياً الجولة 1 من أصل 20 كحد أقصى.} \\
\texttt{لقد جمعتما معاً 1 من أصل 12 عقوبة.} \\
\texttt{أرسل اللاعب الآخر هذه الرسالة:} \\
\\ 
\\ 
\texttt{"مرحباً! لقد استلمت مواقع الكائنات لديك. إليك مواقع الكائنات الموجودة في لوحتي حالياً: W: (3, 6), I: (3, 7), T: (6, 5), C: (3, 4), H: (7, 6). لنحاول الاتفاق على إحداثيات هدف موحدة."} \\
\\ 
\\ 
\texttt{ما هو أمرك التالي؟} \\
            }
        }
    }
    & \\ \\

    \theutterance \stepcounter{utterance}  
    & \multicolumn{4}{p{0.6\linewidth}}{
        \cellcolor[rgb]{0.9,0.9,0.9}{
            \makecell[{{p{\linewidth}}}]{
                \texttt{\tiny{[P1$\rangle$GM]}}
                \texttt{قل: شكراً لك. لنحاول الاتفاق على هدف موحد. ما رأيك أن نحاول نقل جميع الكائنات إلى المواقع التالية: W: (1, 1), I: (2, 1), T: (3, 1), C: (4, 1), H: (5, 1)؟ إذا كنت توافق، سأبدأ بالتحريك.} \\
            }
        }
    }
    & & \\ \\

    \theutterance \stepcounter{utterance}  
    & & & \multicolumn{2}{p{0.3\linewidth}}{
        \cellcolor[rgb]{0.9,0.9,0.9}{
            \makecell[{{p{\linewidth}}}]{
                \texttt{\tiny{[GM$|$GM]}}
                \texttt{قل: شكراً لك. لنحاول الاتفاق على هدف موحد. ما رأيك أن نحاول نقل جميع الكائنات إلى المواقع التالية: W: (1, 1), I: (2, 1), T: (3, 1), C: (4, 1), H: (5, 1)؟ إذا كنت توافق، سأبدأ بالتحريك.} \\
            }
        }
    }
    & & \\ \\

    \theutterance \stepcounter{utterance}  
    & & \multicolumn{4}{p{0.6\linewidth}}{
        \cellcolor[rgb]{0.9,0.9,0.9}{
            \makecell[{{p{\linewidth}}}]{
                \texttt{\tiny{[GM$\rangle$P2]}}
                \texttt{تم تمرير رسالتك إلى اللاعب الآخر.} \\
\\ 
\\ 
\texttt{أنت تلعب حالياً الجولة 2 من أصل 20 كحد أقصى.} \\
\texttt{لقد جمعتما معاً 1 من أصل 12 عقوبة.} \\
\texttt{أرسل اللاعب الآخر هذه الرسالة:} \\
\\ 
\\ 
\texttt{"شكراً لك. لنحاول الاتفاق على هدف موحد. ما رأيك أن نحاول نقل جميع الكائنات إلى المواقع التالية: W: (1, 1), I: (2, 1), T: (3, 1), C: (4, 1), H: (5, 1)؟ إذا كنت توافق، سأبدأ بالتحريك."} \\
\\ 
\\ 
\texttt{ما هو أمرك التالي؟} \\
            }
        }
    }
    & \\ \\

    \theutterance \stepcounter{utterance}  
    & & & \multicolumn{4}{p{0.6\linewidth}}{
        \cellcolor[rgb]{0.9,0.9,0.9}{
            \makecell[{{p{\linewidth}}}]{
                \texttt{\tiny{[GM$\langle$P2]}}
                \texttt{قل: أنا موافق على هذه المواقع المقترحة كهدف مشترك. تفضل بالبدء بالتحريك، وسأقوم أنا أيضاً بتحريك الكائنات من جانبي للوصول إلى هذه الإحداثيات.} \\
            }
        }
    }
     \\ \\

    \theutterance \stepcounter{utterance}  
    & & & \multicolumn{2}{p{0.3\linewidth}}{
        \cellcolor[rgb]{0.9,0.9,0.9}{
            \makecell[{{p{\linewidth}}}]{
                \texttt{\tiny{[GM$|$GM]}}
                \texttt{قل: أنا موافق على هذه المواقع المقترحة كهدف مشترك. تفضل بالبدء بالتحريك، وسأقوم أنا أيضاً بتحريك الكائنات من جانبي للوصول إلى هذه الإحداثيات.} \\
            }
        }
    }
    & & \\ \\

    \theutterance \stepcounter{utterance}  
    & & \multicolumn{4}{p{0.6\linewidth}}{
        \cellcolor[rgb]{0.9,0.9,0.9}{
            \makecell[{{p{\linewidth}}}]{
                \texttt{\tiny{[P1$\langle$GM]}}
                \texttt{تم تمرير رسالتك إلى اللاعب الآخر.} \\
\\ 
\\ 
\texttt{أنت تلعب حالياً الجولة 2 من أصل 20 كحد أقصى.} \\
\texttt{لقد جمعتما معاً 1 من أصل 12 عقوبة.} \\
\texttt{أرسل اللاعب الآخر هذه الرسالة:} \\
\\ 
\\ 
\texttt{"أنا موافق على هذه المواقع المقترحة كهدف مشترك. تفضل بالبدء بالتحريك، وسأقوم أنا أيضاً بتحريك الكائنات من جانبي للوصول إلى هذه الإحداثيات."} \\
\\ 
\\ 
\texttt{ما هو أمرك التالي؟} \\
            }
        }
    }
    & \\ \\

    \theutterance \stepcounter{utterance}  
    & \multicolumn{4}{p{0.6\linewidth}}{
        \cellcolor[rgb]{0.9,0.9,0.9}{
            \makecell[{{p{\linewidth}}}]{
                \texttt{\tiny{[P1$\rangle$GM]}}
                \texttt{حرّك: W, (1, 1)} \\
            }
        }
    }
    & & \\ \\

    \theutterance \stepcounter{utterance}  
    & & & \multicolumn{2}{p{0.3\linewidth}}{
        \cellcolor[rgb]{0.9,0.9,0.9}{
            \makecell[{{p{\linewidth}}}]{
                \texttt{\tiny{[GM$|$GM]}}
                \texttt{حرّك: W, (1, 1)} \\
            }
        }
    }
    & & \\ \\

    \theutterance \stepcounter{utterance}  
    & & & \multicolumn{2}{p{0.3\linewidth}}{
        \cellcolor[rgb]{0.9,0.9,0.9}{
            \makecell[{{p{\linewidth}}}]{
                \texttt{\tiny{[GM$|$GM]}}
                \texttt{W moved to (1, 1)} \\
            }
        }
    }
    & & \\ \\

    \theutterance \stepcounter{utterance}  
    & & \multicolumn{4}{p{0.6\linewidth}}{
        \cellcolor[rgb]{0.9,0.9,0.9}{
            \makecell[{{p{\linewidth}}}]{
                \texttt{\tiny{[GM$\rangle$P2]}}
                \texttt{تم تمرير رسالتك إلى اللاعب الآخر.} \\
\\ 
\\ 
\texttt{أنت تلعب حالياً الجولة 3 من أصل 20 كحد أقصى.} \\
\texttt{لقد جمعتما معاً 1 من أصل 12 عقوبة.} \\
\texttt{حرّك اللاعب الآخر كائناً على شبكته.} \\
\\ 
\\ 
\texttt{ما هو أمرك التالي؟} \\
            }
        }
    }
    & \\ \\

    \theutterance \stepcounter{utterance}  
    & & & \multicolumn{4}{p{0.6\linewidth}}{
        \cellcolor[rgb]{0.9,0.9,0.9}{
            \makecell[{{p{\linewidth}}}]{
                \texttt{\tiny{[GM$\langle$P2]}}
                \texttt{بما أن اللاعب الآخر بدأ بالتحريك، يجب عليّ البدء بتحريك أحد الكائنات من لوحتي نحو الإحداثيات المتفق عليها (W: (1, 1), I: (2, 1), T: (3, 1), C: (4, 1), H: (5, 1)).} \\
\\ 
\texttt{سأبدأ بتحريك الكائن 'C' من موقعه الحالي (3, 4) إلى موضع وسيط أو مباشرة نحو الهدف إذا كان المسار يسمح، ولكن بما أنني لا أستطيع رؤية المسار الكامل، سأحاول تحريك 'C' إلى موضع أقرب للهدف (4, 1).} \\
\\ 
\texttt{**ملاحظة:** يجب التأكد من أن (4, 1) فارغ في لوحتي.} \\
\texttt{بالنظر إلى لوحتي:} \\
\texttt{(4, 1) هو `◌` (فارغ).} \\
\texttt{(3, 4) هو `C`.} \\
\\ 
\texttt{حرّك: C, (4, 1)} \\
            }
        }
    }
     \\ \\

    \theutterance \stepcounter{utterance}  
    & & & \multicolumn{2}{p{0.3\linewidth}}{
        \cellcolor[rgb]{0.9,0.9,0.9}{
            \makecell[{{p{\linewidth}}}]{
                \texttt{\tiny{[GM$|$GM]}}
                \texttt{بما أن اللاعب الآخر بدأ بالتحريك، يجب عليّ البدء بتحريك أحد الكائنات من لوحتي نحو الإحداثيات المتفق عليها (W: (1, 1), I: (2, 1), T: (3, 1), C: (4, 1), H: (5, 1)).} \\
\\ 
\texttt{سأبدأ بتحريك الكائن 'C' من موقعه الحالي (3, 4) إلى موضع وسيط أو مباشرة نحو الهدف إذا كان المسار يسمح، ولكن بما أنني لا أستطيع رؤية المسار الكامل، سأحاول تحريك 'C' إلى موضع أقرب للهدف (4, 1).} \\
\\ 
\texttt{**ملاحظة:** يجب التأكد من أن (4, 1) فارغ في لوحتي.} \\
\texttt{بالنظر إلى لوحتي:} \\
\texttt{(4, 1) هو `◌` (فارغ).} \\
\texttt{(3, 4) هو `C`.} \\
\\ 
\texttt{حرّك: C, (4, 1)} \\
            }
        }
    }
    & & \\ \\

    \theutterance \stepcounter{utterance}  
    & & & \multicolumn{2}{p{0.3\linewidth}}{
        \cellcolor[rgb]{0.9,0.9,0.9}{
            \makecell[{{p{\linewidth}}}]{
                \texttt{\tiny{[GM$|$GM]}}
                \texttt{Invalid format: head is not empty: 'بما أن اللاعب الآخر بدأ بالتحريك، يجب عليّ البدء بتحريك أحد الكائنات من لوحتي نحو الإحداثيات المتفق عليها (W: (1, 1), I: (2, 1), T: (3, 1), C: (4, 1), H: (5, 1)).} \\
\\ 
\texttt{سأبدأ بتحريك الكائن 'C' من موقعه الحالي (3, 4) إلى موضع وسيط أو مباشرة نحو الهدف إذا كان المسار يسمح، ولكن بما أنني لا أستطيع رؤية المسار الكامل، سأحاول تحريك 'C' إلى موضع أقرب للهدف (4, 1).} \\
\\ 
\texttt{**ملاحظة:** يجب التأكد من أن (4, 1) فارغ في لوحتي.} \\
\texttt{بالنظر إلى لوحتي:} \\
\texttt{(4, 1) هو ◌ (فارغ).} \\
\texttt{(3, 4) هو C.} \\
\\ 
\texttt{'} \\
            }
        }
    }
    & & \\ \\

    \theutterance \stepcounter{utterance}  
    & & \multicolumn{4}{p{0.6\linewidth}}{
        \cellcolor[rgb]{0.9,0.9,0.9}{
            \makecell[{{p{\linewidth}}}]{
                \texttt{\tiny{[GM$\rangle$P2]}}
                \texttt{عقوبة: يجب ألا تحتوي رسالتك على أي شيء قبل الأمر!} \\
\\ 
\texttt{تأكد من أن ردك يحتوي فقط على `قل: <رسالة>` أو `حرّك: <كائن>, (<X>, <Y>)`، ولا شيء غير ذلك!} \\
\\ 
\texttt{لقد جمعتما معاً 2 من أصل 12 عقوبة.} \\
\texttt{يرجى المحاولة مرة أخرى!} \\
            }
        }
    }
    & \\ \\

    \theutterance \stepcounter{utterance}  
    & & & \multicolumn{4}{p{0.6\linewidth}}{
        \cellcolor[rgb]{0.9,0.9,0.9}{
            \makecell[{{p{\linewidth}}}]{
                \texttt{\tiny{[GM$\langle$P2]}}
                \texttt{حرّك: C, (4, 1)} \\
            }
        }
    }
     \\ \\

    \theutterance \stepcounter{utterance}  
    & & & \multicolumn{2}{p{0.3\linewidth}}{
        \cellcolor[rgb]{0.9,0.9,0.9}{
            \makecell[{{p{\linewidth}}}]{
                \texttt{\tiny{[GM$|$GM]}}
                \texttt{حرّك: C, (4, 1)} \\
            }
        }
    }
    & & \\ \\

    \theutterance \stepcounter{utterance}  
    & & & \multicolumn{2}{p{0.3\linewidth}}{
        \cellcolor[rgb]{0.9,0.9,0.9}{
            \makecell[{{p{\linewidth}}}]{
                \texttt{\tiny{[GM$|$GM]}}
                \texttt{عقوبة: (4,1) ليست فارغة، بل تحتوي على '│'.} \\
\\ 
\texttt{لقد جمعتما معاً 3 من أصل 12 عقوبة.} \\
\texttt{يرجى المحاولة مرة أخرى!} \\
            }
        }
    }
    & & \\ \\

    \theutterance \stepcounter{utterance}  
    & & \multicolumn{4}{p{0.6\linewidth}}{
        \cellcolor[rgb]{0.9,0.9,0.9}{
            \makecell[{{p{\linewidth}}}]{
                \texttt{\tiny{[GM$\rangle$P2]}}
                \texttt{عقوبة: (4,1) ليست فارغة، بل تحتوي على '│'.} \\
\\ 
\texttt{لقد جمعتما معاً 3 من أصل 12 عقوبة.} \\
\texttt{يرجى المحاولة مرة أخرى!} \\
            }
        }
    }
    & \\ \\

    \theutterance \stepcounter{utterance}  
    & & & \multicolumn{4}{p{0.6\linewidth}}{
        \cellcolor[rgb]{0.9,0.9,0.9}{
            \makecell[{{p{\linewidth}}}]{
                \texttt{\tiny{[GM$\langle$P2]}}
                \texttt{حرّك: C, (4, 2)} \\
            }
        }
    }
     \\ \\

    \theutterance \stepcounter{utterance}  
    & & & \multicolumn{2}{p{0.3\linewidth}}{
        \cellcolor[rgb]{0.9,0.9,0.9}{
            \makecell[{{p{\linewidth}}}]{
                \texttt{\tiny{[GM$|$GM]}}
                \texttt{حرّك: C, (4, 2)} \\
            }
        }
    }
    & & \\ \\

    \theutterance \stepcounter{utterance}  
    & & & \multicolumn{2}{p{0.3\linewidth}}{
        \cellcolor[rgb]{0.9,0.9,0.9}{
            \makecell[{{p{\linewidth}}}]{
                \texttt{\tiny{[GM$|$GM]}}
                \texttt{عقوبة: (4,2) ليست فارغة، بل تحتوي على '┼'.} \\
\\ 
\texttt{لقد جمعتما معاً 4 من أصل 12 عقوبة.} \\
\texttt{يرجى المحاولة مرة أخرى!} \\
            }
        }
    }
    & & \\ \\

    \theutterance \stepcounter{utterance}  
    & & \multicolumn{4}{p{0.6\linewidth}}{
        \cellcolor[rgb]{0.9,0.9,0.9}{
            \makecell[{{p{\linewidth}}}]{
                \texttt{\tiny{[GM$\rangle$P2]}}
                \texttt{عقوبة: (4,2) ليست فارغة، بل تحتوي على '┼'.} \\
\\ 
\texttt{لقد جمعتما معاً 4 من أصل 12 عقوبة.} \\
\texttt{يرجى المحاولة مرة أخرى!} \\
            }
        }
    }
    & \\ \\

    \theutterance \stepcounter{utterance}  
    & & & \multicolumn{4}{p{0.6\linewidth}}{
        \cellcolor[rgb]{0.9,0.9,0.9}{
            \makecell[{{p{\linewidth}}}]{
                \texttt{\tiny{[GM$\langle$P2]}}
                \texttt{حرّك: C, (1, 4)} \\
            }
        }
    }
     \\ \\

    \theutterance \stepcounter{utterance}  
    & & & \multicolumn{2}{p{0.3\linewidth}}{
        \cellcolor[rgb]{0.9,0.9,0.9}{
            \makecell[{{p{\linewidth}}}]{
                \texttt{\tiny{[GM$|$GM]}}
                \texttt{حرّك: C, (1, 4)} \\
            }
        }
    }
    & & \\ \\

    \theutterance \stepcounter{utterance}  
    & & & \multicolumn{2}{p{0.3\linewidth}}{
        \cellcolor[rgb]{0.9,0.9,0.9}{
            \makecell[{{p{\linewidth}}}]{
                \texttt{\tiny{[GM$|$GM]}}
                \texttt{C moved to (1, 4)} \\
            }
        }
    }
    & & \\ \\

    \theutterance \stepcounter{utterance}  
    & & \multicolumn{4}{p{0.6\linewidth}}{
        \cellcolor[rgb]{0.9,0.9,0.9}{
            \makecell[{{p{\linewidth}}}]{
                \texttt{\tiny{[P1$\langle$GM]}}
                \texttt{تم تحريك 'W' إلى (1,1) بنجاح. تبدو شبكتك المحدّثة هكذا:} \\
\texttt{```} \\
\texttt{ 1234567} \\
\texttt{╔═╤═╤═══╗} \\
\texttt{║W│◌│◌◌◌║ 1} \\
\texttt{╟─┴─┼─┬─╢ 2} \\
\texttt{║◌◌◌│◌│T║ 3} \\
\texttt{║◌H◌├─┘◌║ 4} \\
\texttt{║◌┌─┘◌◌◌║ 5} \\
\texttt{╟─┤◌◌◌◌◌║ 6} \\
\texttt{║I│◌C◌◌◌║ 7} \\
\texttt{╚═╧═════╝} \\
\texttt{```} \\
\texttt{أنت تلعب حالياً الجولة 3 من أصل 20 كحد أقصى.} \\
\texttt{لقد جمعتما معاً 4 من أصل 12 عقوبة.} \\
\texttt{حرّك اللاعب الآخر كائناً على شبكته.} \\
\\ 
\\ 
\texttt{ما هو أمرك التالي؟} \\
            }
        }
    }
    & \\ \\

    \theutterance \stepcounter{utterance}  
    & \multicolumn{4}{p{0.6\linewidth}}{
        \cellcolor[rgb]{0.9,0.9,0.9}{
            \makecell[{{p{\linewidth}}}]{
                \texttt{\tiny{[P1$\rangle$GM]}}
                \texttt{حرّك: I, (2, 1)} \\
            }
        }
    }
    & & \\ \\

    \theutterance \stepcounter{utterance}  
    & & & \multicolumn{2}{p{0.3\linewidth}}{
        \cellcolor[rgb]{0.9,0.9,0.9}{
            \makecell[{{p{\linewidth}}}]{
                \texttt{\tiny{[GM$|$GM]}}
                \texttt{حرّك: I, (2, 1)} \\
            }
        }
    }
    & & \\ \\

    \theutterance \stepcounter{utterance}  
    & & & \multicolumn{2}{p{0.3\linewidth}}{
        \cellcolor[rgb]{0.9,0.9,0.9}{
            \makecell[{{p{\linewidth}}}]{
                \texttt{\tiny{[GM$|$GM]}}
                \texttt{عقوبة: (2,1) ليست فارغة، بل تحتوي على '│'.} \\
\\ 
\texttt{لقد جمعتما معاً 5 من أصل 12 عقوبة.} \\
\texttt{يرجى المحاولة مرة أخرى!} \\
            }
        }
    }
    & & \\ \\

    \theutterance \stepcounter{utterance}  
    & & \multicolumn{4}{p{0.6\linewidth}}{
        \cellcolor[rgb]{0.9,0.9,0.9}{
            \makecell[{{p{\linewidth}}}]{
                \texttt{\tiny{[P1$\langle$GM]}}
                \texttt{عقوبة: (2,1) ليست فارغة، بل تحتوي على '│'.} \\
\\ 
\texttt{لقد جمعتما معاً 5 من أصل 12 عقوبة.} \\
\texttt{يرجى المحاولة مرة أخرى!} \\
            }
        }
    }
    & \\ \\

    \theutterance \stepcounter{utterance}  
    & \multicolumn{4}{p{0.6\linewidth}}{
        \cellcolor[rgb]{0.9,0.9,0.9}{
            \makecell[{{p{\linewidth}}}]{
                \texttt{\tiny{[P1$\rangle$GM]}}
                \texttt{حرّك: I, (1, 2)} \\
            }
        }
    }
    & & \\ \\

    \theutterance \stepcounter{utterance}  
    & & & \multicolumn{2}{p{0.3\linewidth}}{
        \cellcolor[rgb]{0.9,0.9,0.9}{
            \makecell[{{p{\linewidth}}}]{
                \texttt{\tiny{[GM$|$GM]}}
                \texttt{حرّك: I, (1, 2)} \\
            }
        }
    }
    & & \\ \\

    \theutterance \stepcounter{utterance}  
    & & & \multicolumn{2}{p{0.3\linewidth}}{
        \cellcolor[rgb]{0.9,0.9,0.9}{
            \makecell[{{p{\linewidth}}}]{
                \texttt{\tiny{[GM$|$GM]}}
                \texttt{عقوبة: (1,2) ليست فارغة، بل تحتوي على '─'.} \\
\\ 
\texttt{لقد جمعتما معاً 6 من أصل 12 عقوبة.} \\
\texttt{يرجى المحاولة مرة أخرى!} \\
            }
        }
    }
    & & \\ \\

    \theutterance \stepcounter{utterance}  
    & & \multicolumn{4}{p{0.6\linewidth}}{
        \cellcolor[rgb]{0.9,0.9,0.9}{
            \makecell[{{p{\linewidth}}}]{
                \texttt{\tiny{[P1$\langle$GM]}}
                \texttt{عقوبة: (1,2) ليست فارغة، بل تحتوي على '─'.} \\
\\ 
\texttt{لقد جمعتما معاً 6 من أصل 12 عقوبة.} \\
\texttt{يرجى المحاولة مرة أخرى!} \\
            }
        }
    }
    & \\ \\

    \theutterance \stepcounter{utterance}  
    & \multicolumn{4}{p{0.6\linewidth}}{
        \cellcolor[rgb]{0.9,0.9,0.9}{
            \makecell[{{p{\linewidth}}}]{
                \texttt{\tiny{[P1$\rangle$GM]}}
                \texttt{حرّك: I, (1, 1)} \\
\\ 
\texttt{(ملاحظة للمدير: بما أن W انتقلت إلى (1,1) في الجولة السابقة، سأحاول تحريك I إلى مكان فارغ متاح. سأحاول تحريك I إلى (1,3) بدلاً من ذلك لتجنب الخطوط والإطارات).} \\
\\ 
\texttt{حرّك: I, (1, 3)} \\
            }
        }
    }
    & & \\ \\

    \theutterance \stepcounter{utterance}  
    & & & \multicolumn{2}{p{0.3\linewidth}}{
        \cellcolor[rgb]{0.9,0.9,0.9}{
            \makecell[{{p{\linewidth}}}]{
                \texttt{\tiny{[GM$|$GM]}}
                \texttt{حرّك: I, (1, 1)} \\
\\ 
\texttt{(ملاحظة للمدير: بما أن W انتقلت إلى (1,1) في الجولة السابقة، سأحاول تحريك I إلى مكان فارغ متاح. سأحاول تحريك I إلى (1,3) بدلاً من ذلك لتجنب الخطوط والإطارات).} \\
\\ 
\texttt{حرّك: I, (1, 3)} \\
            }
        }
    }
    & & \\ \\

    \theutterance \stepcounter{utterance}  
    & & & \multicolumn{2}{p{0.3\linewidth}}{
        \cellcolor[rgb]{0.9,0.9,0.9}{
            \makecell[{{p{\linewidth}}}]{
                \texttt{\tiny{[GM$|$GM]}}
                \texttt{Invalid format: head is not empty: 'حرّك: I, (1, 1)} \\
\\ 
\texttt{(ملاحظة للمدير: بما أن W انتقلت إلى (1,1) في الجولة السابقة، سأحاول تحريك I إلى مكان فارغ متاح. سأحاول تحريك I إلى (1,3) بدلاً من ذلك لتجنب الخطوط والإطارات).} \\
\\ 
\texttt{'} \\
            }
        }
    }
    & & \\ \\

    \theutterance \stepcounter{utterance}  
    & & \multicolumn{4}{p{0.6\linewidth}}{
        \cellcolor[rgb]{0.9,0.9,0.9}{
            \makecell[{{p{\linewidth}}}]{
                \texttt{\tiny{[P1$\langle$GM]}}
                \texttt{عقوبة: يجب ألا تحتوي رسالتك على أي شيء قبل الأمر!} \\
\\ 
\texttt{تأكد من أن ردك يحتوي فقط على `قل: <رسالة>` أو `حرّك: <كائن>, (<X>, <Y>)`، ولا شيء غير ذلك!} \\
\\ 
\texttt{لقد جمعتما معاً 7 من أصل 12 عقوبة.} \\
\texttt{يرجى المحاولة مرة أخرى!} \\
            }
        }
    }
    & \\ \\

    \theutterance \stepcounter{utterance}  
    & \multicolumn{4}{p{0.6\linewidth}}{
        \cellcolor[rgb]{0.9,0.9,0.9}{
            \makecell[{{p{\linewidth}}}]{
                \texttt{\tiny{[P1$\rangle$GM]}}
                \texttt{حرّك: I, (1, 3)} \\
            }
        }
    }
    & & \\ \\

    \theutterance \stepcounter{utterance}  
    & & & \multicolumn{2}{p{0.3\linewidth}}{
        \cellcolor[rgb]{0.9,0.9,0.9}{
            \makecell[{{p{\linewidth}}}]{
                \texttt{\tiny{[GM$|$GM]}}
                \texttt{حرّك: I, (1, 3)} \\
            }
        }
    }
    & & \\ \\

    \theutterance \stepcounter{utterance}  
    & & & \multicolumn{2}{p{0.3\linewidth}}{
        \cellcolor[rgb]{0.9,0.9,0.9}{
            \makecell[{{p{\linewidth}}}]{
                \texttt{\tiny{[GM$|$GM]}}
                \texttt{I moved to (1, 3)} \\
            }
        }
    }
    & & \\ \\

    \theutterance \stepcounter{utterance}  
    & & \multicolumn{4}{p{0.6\linewidth}}{
        \cellcolor[rgb]{0.9,0.9,0.9}{
            \makecell[{{p{\linewidth}}}]{
                \texttt{\tiny{[GM$\rangle$P2]}}
                \texttt{تم تحريك 'C' إلى (1,4) بنجاح. تبدو شبكتك المحدّثة هكذا:} \\
\texttt{```} \\
\texttt{ 1234567} \\
\texttt{╔═╤═╤═══╗} \\
\texttt{║◌│◌│◌◌◌║ 1} \\
\texttt{╟─┴─┼─┬─╢ 2} \\
\texttt{║◌◌◌│◌│◌║ 3} \\
\texttt{║C◌◌├─┘◌║ 4} \\
\texttt{║◌┌─┘◌T◌║ 5} \\
\texttt{╟─┤◌W◌◌H║ 6} \\
\texttt{║◌│◌I◌◌◌║ 7} \\
\texttt{╚═╧═════╝} \\
\texttt{```} \\
\texttt{أنت تلعب حالياً الجولة 4 من أصل 20 كحد أقصى.} \\
\texttt{لقد جمعتما معاً 7 من أصل 12 عقوبة.} \\
\texttt{حرّك اللاعب الآخر كائناً على شبكته.} \\
\\ 
\\ 
\texttt{ما هو أمرك التالي؟} \\
            }
        }
    }
    & \\ \\

    \theutterance \stepcounter{utterance}  
    & & & \multicolumn{4}{p{0.6\linewidth}}{
        \cellcolor[rgb]{0.9,0.9,0.9}{
            \makecell[{{p{\linewidth}}}]{
                \texttt{\tiny{[GM$\langle$P2]}}
                \texttt{حرّك: C, (1, 1)} \\
            }
        }
    }
     \\ \\

    \theutterance \stepcounter{utterance}  
    & & & \multicolumn{2}{p{0.3\linewidth}}{
        \cellcolor[rgb]{0.9,0.9,0.9}{
            \makecell[{{p{\linewidth}}}]{
                \texttt{\tiny{[GM$|$GM]}}
                \texttt{حرّك: C, (1, 1)} \\
            }
        }
    }
    & & \\ \\

    \theutterance \stepcounter{utterance}  
    & & & \multicolumn{2}{p{0.3\linewidth}}{
        \cellcolor[rgb]{0.9,0.9,0.9}{
            \makecell[{{p{\linewidth}}}]{
                \texttt{\tiny{[GM$|$GM]}}
                \texttt{C moved to (1, 1)} \\
            }
        }
    }
    & & \\ \\

    \theutterance \stepcounter{utterance}  
    & & \multicolumn{4}{p{0.6\linewidth}}{
        \cellcolor[rgb]{0.9,0.9,0.9}{
            \makecell[{{p{\linewidth}}}]{
                \texttt{\tiny{[P1$\langle$GM]}}
                \texttt{تم تحريك 'I' إلى (1,3) بنجاح. تبدو شبكتك المحدّثة هكذا:} \\
\texttt{```} \\
\texttt{ 1234567} \\
\texttt{╔═╤═╤═══╗} \\
\texttt{║W│◌│◌◌◌║ 1} \\
\texttt{╟─┴─┼─┬─╢ 2} \\
\texttt{║I◌◌│◌│T║ 3} \\
\texttt{║◌H◌├─┘◌║ 4} \\
\texttt{║◌┌─┘◌◌◌║ 5} \\
\texttt{╟─┤◌◌◌◌◌║ 6} \\
\texttt{║◌│◌C◌◌◌║ 7} \\
\texttt{╚═╧═════╝} \\
\texttt{```} \\
\texttt{أنت تلعب حالياً الجولة 4 من أصل 20 كحد أقصى.} \\
\texttt{لقد جمعتما معاً 7 من أصل 12 عقوبة.} \\
\texttt{حرّك اللاعب الآخر كائناً على شبكته.} \\
\\ 
\\ 
\texttt{ما هو أمرك التالي؟} \\
            }
        }
    }
    & \\ \\

    \theutterance \stepcounter{utterance}  
    & \multicolumn{4}{p{0.6\linewidth}}{
        \cellcolor[rgb]{0.9,0.9,0.9}{
            \makecell[{{p{\linewidth}}}]{
                \texttt{\tiny{[P1$\rangle$GM]}}
                \texttt{حرّك: T, (1, 1)} \\
\\ 
\texttt{(ملاحظة: سأحاول تحريك T إلى (1,1) ولكن W موجودة هناك. سأختار مكاناً فارغاً بدلاً من ذلك. سأحاول تحريك T إلى (5, 1) حيث يبدو فارغاً).} \\
\\ 
\texttt{حرّك: T, (5, 1)} \\
            }
        }
    }
    & & \\ \\

    \theutterance \stepcounter{utterance}  
    & & & \multicolumn{2}{p{0.3\linewidth}}{
        \cellcolor[rgb]{0.9,0.9,0.9}{
            \makecell[{{p{\linewidth}}}]{
                \texttt{\tiny{[GM$|$GM]}}
                \texttt{حرّك: T, (1, 1)} \\
\\ 
\texttt{(ملاحظة: سأحاول تحريك T إلى (1,1) ولكن W موجودة هناك. سأختار مكاناً فارغاً بدلاً من ذلك. سأحاول تحريك T إلى (5, 1) حيث يبدو فارغاً).} \\
\\ 
\texttt{حرّك: T, (5, 1)} \\
            }
        }
    }
    & & \\ \\

    \theutterance \stepcounter{utterance}  
    & & & \multicolumn{2}{p{0.3\linewidth}}{
        \cellcolor[rgb]{0.9,0.9,0.9}{
            \makecell[{{p{\linewidth}}}]{
                \texttt{\tiny{[GM$|$GM]}}
                \texttt{Invalid format: head is not empty: 'حرّك: T, (1, 1)} \\
\\ 
\texttt{(ملاحظة: سأحاول تحريك T إلى (1,1) ولكن W موجودة هناك. سأختار مكاناً فارغاً بدلاً من ذلك. سأحاول تحريك T إلى (5, 1) حيث يبدو فارغاً).} \\
\\ 
\texttt{'} \\
            }
        }
    }
    & & \\ \\

    \theutterance \stepcounter{utterance}  
    & & \multicolumn{4}{p{0.6\linewidth}}{
        \cellcolor[rgb]{0.9,0.9,0.9}{
            \makecell[{{p{\linewidth}}}]{
                \texttt{\tiny{[P1$\langle$GM]}}
                \texttt{عقوبة: يجب ألا تحتوي رسالتك على أي شيء قبل الأمر!} \\
\\ 
\texttt{تأكد من أن ردك يحتوي فقط على `قل: <رسالة>` أو `حرّك: <كائن>, (<X>, <Y>)`، ولا شيء غير ذلك!} \\
\\ 
\texttt{لقد جمعتما معاً 8 من أصل 12 عقوبة.} \\
\texttt{يرجى المحاولة مرة أخرى!} \\
            }
        }
    }
    & \\ \\

    \theutterance \stepcounter{utterance}  
    & \multicolumn{4}{p{0.6\linewidth}}{
        \cellcolor[rgb]{0.9,0.9,0.9}{
            \makecell[{{p{\linewidth}}}]{
                \texttt{\tiny{[P1$\rangle$GM]}}
                \texttt{حرّك: T, (5, 1)} \\
            }
        }
    }
    & & \\ \\

    \theutterance \stepcounter{utterance}  
    & & & \multicolumn{2}{p{0.3\linewidth}}{
        \cellcolor[rgb]{0.9,0.9,0.9}{
            \makecell[{{p{\linewidth}}}]{
                \texttt{\tiny{[GM$|$GM]}}
                \texttt{حرّك: T, (5, 1)} \\
            }
        }
    }
    & & \\ \\

    \theutterance \stepcounter{utterance}  
    & & & \multicolumn{2}{p{0.3\linewidth}}{
        \cellcolor[rgb]{0.9,0.9,0.9}{
            \makecell[{{p{\linewidth}}}]{
                \texttt{\tiny{[GM$|$GM]}}
                \texttt{T moved to (5, 1)} \\
            }
        }
    }
    & & \\ \\

    \theutterance \stepcounter{utterance}  
    & & \multicolumn{4}{p{0.6\linewidth}}{
        \cellcolor[rgb]{0.9,0.9,0.9}{
            \makecell[{{p{\linewidth}}}]{
                \texttt{\tiny{[GM$\rangle$P2]}}
                \texttt{تم تحريك 'C' إلى (1,1) بنجاح. تبدو شبكتك المحدّثة هكذا:} \\
\texttt{```} \\
\texttt{ 1234567} \\
\texttt{╔═╤═╤═══╗} \\
\texttt{║C│◌│◌◌◌║ 1} \\
\texttt{╟─┴─┼─┬─╢ 2} \\
\texttt{║◌◌◌│◌│◌║ 3} \\
\texttt{║◌◌◌├─┘◌║ 4} \\
\texttt{║◌┌─┘◌T◌║ 5} \\
\texttt{╟─┤◌W◌◌H║ 6} \\
\texttt{║◌│◌I◌◌◌║ 7} \\
\texttt{╚═╧═════╝} \\
\texttt{```} \\
\texttt{أنت تلعب حالياً الجولة 5 من أصل 20 كحد أقصى.} \\
\texttt{لقد جمعتما معاً 8 من أصل 12 عقوبة.} \\
\texttt{حرّك اللاعب الآخر كائناً على شبكته.} \\
\\ 
\\ 
\texttt{ما هو أمرك التالي؟} \\
            }
        }
    }
    & \\ \\

    \theutterance \stepcounter{utterance}  
    & & & \multicolumn{4}{p{0.6\linewidth}}{
        \cellcolor[rgb]{0.9,0.9,0.9}{
            \makecell[{{p{\linewidth}}}]{
                \texttt{\tiny{[GM$\langle$P2]}}
                \texttt{حرّك: I, (2, 1)} \\
            }
        }
    }
     \\ \\

    \theutterance \stepcounter{utterance}  
    & & & \multicolumn{2}{p{0.3\linewidth}}{
        \cellcolor[rgb]{0.9,0.9,0.9}{
            \makecell[{{p{\linewidth}}}]{
                \texttt{\tiny{[GM$|$GM]}}
                \texttt{حرّك: I, (2, 1)} \\
            }
        }
    }
    & & \\ \\

    \theutterance \stepcounter{utterance}  
    & & & \multicolumn{2}{p{0.3\linewidth}}{
        \cellcolor[rgb]{0.9,0.9,0.9}{
            \makecell[{{p{\linewidth}}}]{
                \texttt{\tiny{[GM$|$GM]}}
                \texttt{عقوبة: (2,1) ليست فارغة، بل تحتوي على '│'.} \\
\\ 
\texttt{لقد جمعتما معاً 9 من أصل 12 عقوبة.} \\
\texttt{يرجى المحاولة مرة أخرى!} \\
            }
        }
    }
    & & \\ \\

    \theutterance \stepcounter{utterance}  
    & & \multicolumn{4}{p{0.6\linewidth}}{
        \cellcolor[rgb]{0.9,0.9,0.9}{
            \makecell[{{p{\linewidth}}}]{
                \texttt{\tiny{[GM$\rangle$P2]}}
                \texttt{عقوبة: (2,1) ليست فارغة، بل تحتوي على '│'.} \\
\\ 
\texttt{لقد جمعتما معاً 9 من أصل 12 عقوبة.} \\
\texttt{يرجى المحاولة مرة أخرى!} \\
            }
        }
    }
    & \\ \\

    \theutterance \stepcounter{utterance}  
    & & & \multicolumn{4}{p{0.6\linewidth}}{
        \cellcolor[rgb]{0.9,0.9,0.9}{
            \makecell[{{p{\linewidth}}}]{
                \texttt{\tiny{[GM$\langle$P2]}}
                \texttt{حرّك: I, (3, 1)} \\
            }
        }
    }
     \\ \\

    \theutterance \stepcounter{utterance}  
    & & & \multicolumn{2}{p{0.3\linewidth}}{
        \cellcolor[rgb]{0.9,0.9,0.9}{
            \makecell[{{p{\linewidth}}}]{
                \texttt{\tiny{[GM$|$GM]}}
                \texttt{حرّك: I, (3, 1)} \\
            }
        }
    }
    & & \\ \\

    \theutterance \stepcounter{utterance}  
    & & & \multicolumn{2}{p{0.3\linewidth}}{
        \cellcolor[rgb]{0.9,0.9,0.9}{
            \makecell[{{p{\linewidth}}}]{
                \texttt{\tiny{[GM$|$GM]}}
                \texttt{I moved to (3, 1)} \\
            }
        }
    }
    & & \\ \\

    \theutterance \stepcounter{utterance}  
    & & \multicolumn{4}{p{0.6\linewidth}}{
        \cellcolor[rgb]{0.9,0.9,0.9}{
            \makecell[{{p{\linewidth}}}]{
                \texttt{\tiny{[P1$\langle$GM]}}
                \texttt{تم تحريك 'T' إلى (5,1) بنجاح. تبدو شبكتك المحدّثة هكذا:} \\
\texttt{```} \\
\texttt{ 1234567} \\
\texttt{╔═╤═╤═══╗} \\
\texttt{║W│◌│T◌◌║ 1} \\
\texttt{╟─┴─┼─┬─╢ 2} \\
\texttt{║I◌◌│◌│◌║ 3} \\
\texttt{║◌H◌├─┘◌║ 4} \\
\texttt{║◌┌─┘◌◌◌║ 5} \\
\texttt{╟─┤◌◌◌◌◌║ 6} \\
\texttt{║◌│◌C◌◌◌║ 7} \\
\texttt{╚═╧═════╝} \\
\texttt{```} \\
\texttt{أنت تلعب حالياً الجولة 5 من أصل 20 كحد أقصى.} \\
\texttt{لقد جمعتما معاً 9 من أصل 12 عقوبة.} \\
\texttt{حرّك اللاعب الآخر كائناً على شبكته.} \\
\\ 
\\ 
\texttt{ما هو أمرك التالي؟} \\
            }
        }
    }
    & \\ \\

    \theutterance \stepcounter{utterance}  
    & \multicolumn{4}{p{0.6\linewidth}}{
        \cellcolor[rgb]{0.9,0.9,0.9}{
            \makecell[{{p{\linewidth}}}]{
                \texttt{\tiny{[P1$\rangle$GM]}}
                \texttt{حرّك: H, (2, 1)} \\
            }
        }
    }
    & & \\ \\

    \theutterance \stepcounter{utterance}  
    & & & \multicolumn{2}{p{0.3\linewidth}}{
        \cellcolor[rgb]{0.9,0.9,0.9}{
            \makecell[{{p{\linewidth}}}]{
                \texttt{\tiny{[GM$|$GM]}}
                \texttt{حرّك: H, (2, 1)} \\
            }
        }
    }
    & & \\ \\

    \theutterance \stepcounter{utterance}  
    & & & \multicolumn{2}{p{0.3\linewidth}}{
        \cellcolor[rgb]{0.9,0.9,0.9}{
            \makecell[{{p{\linewidth}}}]{
                \texttt{\tiny{[GM$|$GM]}}
                \texttt{عقوبة: (2,1) ليست فارغة، بل تحتوي على '│'.} \\
\\ 
\texttt{لقد جمعتما معاً 10 من أصل 12 عقوبة.} \\
\texttt{يرجى المحاولة مرة أخرى!} \\
            }
        }
    }
    & & \\ \\

    \theutterance \stepcounter{utterance}  
    & & \multicolumn{4}{p{0.6\linewidth}}{
        \cellcolor[rgb]{0.9,0.9,0.9}{
            \makecell[{{p{\linewidth}}}]{
                \texttt{\tiny{[P1$\langle$GM]}}
                \texttt{عقوبة: (2,1) ليست فارغة، بل تحتوي على '│'.} \\
\\ 
\texttt{لقد جمعتما معاً 10 من أصل 12 عقوبة.} \\
\texttt{يرجى المحاولة مرة أخرى!} \\
            }
        }
    }
    & \\ \\

    \theutterance \stepcounter{utterance}  
    & \multicolumn{4}{p{0.6\linewidth}}{
        \cellcolor[rgb]{0.9,0.9,0.9}{
            \makecell[{{p{\linewidth}}}]{
                \texttt{\tiny{[P1$\rangle$GM]}}
                \texttt{حرّك: H, (2, 3)} \\
            }
        }
    }
    & & \\ \\

    \theutterance \stepcounter{utterance}  
    & & & \multicolumn{2}{p{0.3\linewidth}}{
        \cellcolor[rgb]{0.9,0.9,0.9}{
            \makecell[{{p{\linewidth}}}]{
                \texttt{\tiny{[GM$|$GM]}}
                \texttt{حرّك: H, (2, 3)} \\
            }
        }
    }
    & & \\ \\

    \theutterance \stepcounter{utterance}  
    & & & \multicolumn{2}{p{0.3\linewidth}}{
        \cellcolor[rgb]{0.9,0.9,0.9}{
            \makecell[{{p{\linewidth}}}]{
                \texttt{\tiny{[GM$|$GM]}}
                \texttt{H moved to (2, 3)} \\
            }
        }
    }
    & & \\ \\

    \theutterance \stepcounter{utterance}  
    & & \multicolumn{4}{p{0.6\linewidth}}{
        \cellcolor[rgb]{0.9,0.9,0.9}{
            \makecell[{{p{\linewidth}}}]{
                \texttt{\tiny{[GM$\rangle$P2]}}
                \texttt{تم تحريك 'I' إلى (3,1) بنجاح. تبدو شبكتك المحدّثة هكذا:} \\
\texttt{```} \\
\texttt{ 1234567} \\
\texttt{╔═╤═╤═══╗} \\
\texttt{║C│I│◌◌◌║ 1} \\
\texttt{╟─┴─┼─┬─╢ 2} \\
\texttt{║◌◌◌│◌│◌║ 3} \\
\texttt{║◌◌◌├─┘◌║ 4} \\
\texttt{║◌┌─┘◌T◌║ 5} \\
\texttt{╟─┤◌W◌◌H║ 6} \\
\texttt{║◌│◌◌◌◌◌║ 7} \\
\texttt{╚═╧═════╝} \\
\texttt{```} \\
\texttt{أنت تلعب حالياً الجولة 6 من أصل 20 كحد أقصى.} \\
\texttt{لقد جمعتما معاً 10 من أصل 12 عقوبة.} \\
\texttt{حرّك اللاعب الآخر كائناً على شبكته.} \\
\\ 
\\ 
\texttt{ما هو أمرك التالي؟} \\
            }
        }
    }
    & \\ \\

    \theutterance \stepcounter{utterance}  
    & & & \multicolumn{4}{p{0.6\linewidth}}{
        \cellcolor[rgb]{0.9,0.9,0.9}{
            \makecell[{{p{\linewidth}}}]{
                \texttt{\tiny{[GM$\langle$P2]}}
                \texttt{حرّك: T, (3, 1)} \\
\\ 
\texttt{(ملاحظة: بما أن 'I' أصبحت في (3,1)، يجب عليّ أولاً تحريك 'I' من مكانها أو تحريك 'T' إلى مكان آخر. سأحاول تحريك 'T' إلى (3, 3) أولاً لتفريغ المسار أو تقريبه من الهدف (3,1) دون التصادم).} \\
\\ 
\texttt{عذراً، سأصحح الأمر:} \\
\texttt{حرّك: T, (3, 3)} \\
            }
        }
    }
     \\ \\

    \theutterance \stepcounter{utterance}  
    & & & \multicolumn{2}{p{0.3\linewidth}}{
        \cellcolor[rgb]{0.9,0.9,0.9}{
            \makecell[{{p{\linewidth}}}]{
                \texttt{\tiny{[GM$|$GM]}}
                \texttt{حرّك: T, (3, 1)} \\
\\ 
\texttt{(ملاحظة: بما أن 'I' أصبحت في (3,1)، يجب عليّ أولاً تحريك 'I' من مكانها أو تحريك 'T' إلى مكان آخر. سأحاول تحريك 'T' إلى (3, 3) أولاً لتفريغ المسار أو تقريبه من الهدف (3,1) دون التصادم).} \\
\\ 
\texttt{عذراً، سأصحح الأمر:} \\
\texttt{حرّك: T, (3, 3)} \\
            }
        }
    }
    & & \\ \\

    \theutterance \stepcounter{utterance}  
    & & & \multicolumn{2}{p{0.3\linewidth}}{
        \cellcolor[rgb]{0.9,0.9,0.9}{
            \makecell[{{p{\linewidth}}}]{
                \texttt{\tiny{[GM$|$GM]}}
                \texttt{Invalid format: head is not empty: 'حرّك: T, (3, 1)} \\
\\ 
\texttt{(ملاحظة: بما أن 'I' أصبحت في (3,1)، يجب عليّ أولاً تحريك 'I' من مكانها أو تحريك 'T' إلى مكان آخر. سأحاول تحريك 'T' إلى (3, 3) أولاً لتفريغ المسار أو تقريبه من الهدف (3,1) دون التصادم).} \\
\\ 
\texttt{عذراً، سأصحح الأمر:} \\
\texttt{'} \\
            }
        }
    }
    & & \\ \\

    \theutterance \stepcounter{utterance}  
    & & \multicolumn{4}{p{0.6\linewidth}}{
        \cellcolor[rgb]{0.9,0.9,0.9}{
            \makecell[{{p{\linewidth}}}]{
                \texttt{\tiny{[GM$\rangle$P2]}}
                \texttt{عقوبة: يجب ألا تحتوي رسالتك على أي شيء قبل الأمر!} \\
\\ 
\texttt{تأكد من أن ردك يحتوي فقط على `قل: <رسالة>` أو `حرّك: <كائن>, (<X>, <Y>)`، ولا شيء غير ذلك!} \\
\\ 
\texttt{لقد جمعتما معاً 11 من أصل 12 عقوبة.} \\
\texttt{يرجى المحاولة مرة أخرى!} \\
            }
        }
    }
    & \\ \\

    \theutterance \stepcounter{utterance}  
    & & & \multicolumn{4}{p{0.6\linewidth}}{
        \cellcolor[rgb]{0.9,0.9,0.9}{
            \makecell[{{p{\linewidth}}}]{
                \texttt{\tiny{[GM$\langle$P2]}}
                \texttt{حرّك: T, (3, 3)} \\
            }
        }
    }
     \\ \\

    \theutterance \stepcounter{utterance}  
    & & & \multicolumn{2}{p{0.3\linewidth}}{
        \cellcolor[rgb]{0.9,0.9,0.9}{
            \makecell[{{p{\linewidth}}}]{
                \texttt{\tiny{[GM$|$GM]}}
                \texttt{حرّك: T, (3, 3)} \\
            }
        }
    }
    & & \\ \\

    \theutterance \stepcounter{utterance}  
    & & & \multicolumn{2}{p{0.3\linewidth}}{
        \cellcolor[rgb]{0.9,0.9,0.9}{
            \makecell[{{p{\linewidth}}}]{
                \texttt{\tiny{[GM$|$GM]}}
                \texttt{T moved to (3, 3)} \\
            }
        }
    }
    & & \\ \\

    \theutterance \stepcounter{utterance}  
    & & \multicolumn{4}{p{0.6\linewidth}}{
        \cellcolor[rgb]{0.9,0.9,0.9}{
            \makecell[{{p{\linewidth}}}]{
                \texttt{\tiny{[P1$\langle$GM]}}
                \texttt{تم تحريك 'H' إلى (2,3) بنجاح. تبدو شبكتك المحدّثة هكذا:} \\
\texttt{```} \\
\texttt{ 1234567} \\
\texttt{╔═╤═╤═══╗} \\
\texttt{║W│◌│T◌◌║ 1} \\
\texttt{╟─┴─┼─┬─╢ 2} \\
\texttt{║IH◌│◌│◌║ 3} \\
\texttt{║◌◌◌├─┘◌║ 4} \\
\texttt{║◌┌─┘◌◌◌║ 5} \\
\texttt{╟─┤◌◌◌◌◌║ 6} \\
\texttt{║◌│◌C◌◌◌║ 7} \\
\texttt{╚═╧═════╝} \\
\texttt{```} \\
\texttt{أنت تلعب حالياً الجولة 6 من أصل 20 كحد أقصى.} \\
\texttt{لقد جمعتما معاً 11 من أصل 12 عقوبة.} \\
\texttt{حرّك اللاعب الآخر كائناً على شبكته.} \\
\\ 
\\ 
\texttt{ما هو أمرك التالي؟} \\
            }
        }
    }
    & \\ \\

    \theutterance \stepcounter{utterance}  
    & \multicolumn{4}{p{0.6\linewidth}}{
        \cellcolor[rgb]{0.9,0.9,0.9}{
            \makecell[{{p{\linewidth}}}]{
                \texttt{\tiny{[P1$\rangle$GM]}}
                \texttt{حرّك: C, (4, 1)} \\
            }
        }
    }
    & & \\ \\

    \theutterance \stepcounter{utterance}  
    & & & \multicolumn{2}{p{0.3\linewidth}}{
        \cellcolor[rgb]{0.9,0.9,0.9}{
            \makecell[{{p{\linewidth}}}]{
                \texttt{\tiny{[GM$|$GM]}}
                \texttt{حرّك: C, (4, 1)} \\
            }
        }
    }
    & & \\ \\

    \theutterance \stepcounter{utterance}  
    & & & \multicolumn{2}{p{0.3\linewidth}}{
        \cellcolor[rgb]{0.9,0.9,0.9}{
            \makecell[{{p{\linewidth}}}]{
                \texttt{\tiny{[GM$|$GM]}}
                \texttt{عقوبة: (4,1) ليست فارغة، بل تحتوي على '│'.} \\
\\ 
\texttt{لقد جمعتما معاً 12 من أصل 12 عقوبة.} \\
\texttt{يرجى المحاولة مرة أخرى!} \\
            }
        }
    }
    & & \\ \\

    \theutterance \stepcounter{utterance}  
    & & \multicolumn{4}{p{0.6\linewidth}}{
        \cellcolor[rgb]{0.9,0.9,0.9}{
            \makecell[{{p{\linewidth}}}]{
                \texttt{\tiny{[P1$\langle$GM]}}
                \texttt{عقوبة: (4,1) ليست فارغة، بل تحتوي على '│'.} \\
\\ 
\texttt{لقد جمعتما معاً 12 من أصل 12 عقوبة.} \\
\texttt{يرجى المحاولة مرة أخرى!} \\
            }
        }
    }
    & \\ \\

    \theutterance \stepcounter{utterance}  
    & \multicolumn{4}{p{0.6\linewidth}}{
        \cellcolor[rgb]{0.9,0.9,0.9}{
            \makecell[{{p{\linewidth}}}]{
                \texttt{\tiny{[P1$\rangle$GM]}}
                \texttt{حرّك: C, (4, 2)} \\
            }
        }
    }
    & & \\ \\

    \theutterance \stepcounter{utterance}  
    & & & \multicolumn{2}{p{0.3\linewidth}}{
        \cellcolor[rgb]{0.9,0.9,0.9}{
            \makecell[{{p{\linewidth}}}]{
                \texttt{\tiny{[GM$|$GM]}}
                \texttt{حرّك: C, (4, 2)} \\
            }
        }
    }
    & & \\ \\

    \theutterance \stepcounter{utterance}  
    & & & \multicolumn{2}{p{0.3\linewidth}}{
        \cellcolor[rgb]{0.9,0.9,0.9}{
            \makecell[{{p{\linewidth}}}]{
                \texttt{\tiny{[GM$|$GM]}}
                \texttt{عقوبة: (4,2) ليست فارغة، بل تحتوي على '┼'.} \\
\\ 
\texttt{لقد جمعتما معاً 13 من أصل 12 عقوبة.} \\
\texttt{يرجى المحاولة مرة أخرى!} \\
            }
        }
    }
    & & \\ \\

    \theutterance \stepcounter{utterance}  
    & & \multicolumn{4}{p{0.6\linewidth}}{
        \cellcolor[rgb]{0.9,0.9,0.9}{
            \makecell[{{p{\linewidth}}}]{
                \texttt{\tiny{[P1$\langle$GM]}}
                \texttt{عقوبة: (4,2) ليست فارغة، بل تحتوي على '┼'.} \\
\\ 
\texttt{لقد جمعتما معاً 13 من أصل 12 عقوبة.} \\
\texttt{يرجى المحاولة مرة أخرى!} \\
            }
        }
    }
    & \\ \\

    \theutterance \stepcounter{utterance}  
    & \multicolumn{4}{p{0.6\linewidth}}{
        \cellcolor[rgb]{0.9,0.9,0.9}{
            \makecell[{{p{\linewidth}}}]{
                \texttt{\tiny{[P1$\rangle$GM]}}
                \texttt{حرّك: C, (4, 3)} \\
            }
        }
    }
    & & \\ \\

    \theutterance \stepcounter{utterance}  
    & & & \multicolumn{2}{p{0.3\linewidth}}{
        \cellcolor[rgb]{0.9,0.9,0.9}{
            \makecell[{{p{\linewidth}}}]{
                \texttt{\tiny{[GM$|$GM]}}
                \texttt{حرّك: C, (4, 3)} \\
            }
        }
    }
    & & \\ \\

    \theutterance \stepcounter{utterance}  
    & & & \multicolumn{2}{p{0.3\linewidth}}{
        \cellcolor[rgb]{0.9,0.9,0.9}{
            \makecell[{{p{\linewidth}}}]{
                \texttt{\tiny{[GM$|$GM]}}
                \texttt{عقوبة: (4,3) ليست فارغة، بل تحتوي على '│'.} \\
\\ 
\texttt{لقد جمعتما معاً 14 من أصل 12 عقوبة.} \\
\texttt{يرجى المحاولة مرة أخرى!} \\
            }
        }
    }
    & & \\ \\

    \theutterance \stepcounter{utterance}  
    & & \multicolumn{4}{p{0.6\linewidth}}{
        \cellcolor[rgb]{0.9,0.9,0.9}{
            \makecell[{{p{\linewidth}}}]{
                \texttt{\tiny{[P1$\langle$GM]}}
                \texttt{عقوبة: (4,3) ليست فارغة، بل تحتوي على '│'.} \\
\\ 
\texttt{لقد جمعتما معاً 14 من أصل 12 عقوبة.} \\
\texttt{يرجى المحاولة مرة أخرى!} \\
            }
        }
    }
    & \\ \\

    \theutterance \stepcounter{utterance}  
    & \multicolumn{4}{p{0.6\linewidth}}{
        \cellcolor[rgb]{0.9,0.9,0.9}{
            \makecell[{{p{\linewidth}}}]{
                \texttt{\tiny{[P1$\rangle$GM]}}
                \texttt{حرّك: C, (4, 5)} \\
            }
        }
    }
    & & \\ \\

    \theutterance \stepcounter{utterance}  
    & & & \multicolumn{2}{p{0.3\linewidth}}{
        \cellcolor[rgb]{0.9,0.9,0.9}{
            \makecell[{{p{\linewidth}}}]{
                \texttt{\tiny{[GM$|$GM]}}
                \texttt{حرّك: C, (4, 5)} \\
            }
        }
    }
    & & \\ \\

    \theutterance \stepcounter{utterance}  
    & & & \multicolumn{2}{p{0.3\linewidth}}{
        \cellcolor[rgb]{0.9,0.9,0.9}{
            \makecell[{{p{\linewidth}}}]{
                \texttt{\tiny{[GM$|$GM]}}
                \texttt{عقوبة: (4,5) ليست فارغة، بل تحتوي على '┘'.} \\
\\ 
\texttt{لقد جمعتما معاً 15 من أصل 12 عقوبة.} \\
\texttt{يرجى المحاولة مرة أخرى!} \\
            }
        }
    }
    & & \\ \\

    \theutterance \stepcounter{utterance}  
    & & \multicolumn{4}{p{0.6\linewidth}}{
        \cellcolor[rgb]{0.9,0.9,0.9}{
            \makecell[{{p{\linewidth}}}]{
                \texttt{\tiny{[P1$\langle$GM]}}
                \texttt{عقوبة: (4,5) ليست فارغة، بل تحتوي على '┘'.} \\
\\ 
\texttt{لقد جمعتما معاً 15 من أصل 12 عقوبة.} \\
\texttt{يرجى المحاولة مرة أخرى!} \\
            }
        }
    }
    & \\ \\

    \theutterance \stepcounter{utterance}  
    & \multicolumn{4}{p{0.6\linewidth}}{
        \cellcolor[rgb]{0.9,0.9,0.9}{
            \makecell[{{p{\linewidth}}}]{
                \texttt{\tiny{[P1$\rangle$GM]}}
                \texttt{حرّك: C, (4, 6)} \\
            }
        }
    }
    & & \\ \\

    \theutterance \stepcounter{utterance}  
    & & & \multicolumn{2}{p{0.3\linewidth}}{
        \cellcolor[rgb]{0.9,0.9,0.9}{
            \makecell[{{p{\linewidth}}}]{
                \texttt{\tiny{[GM$|$GM]}}
                \texttt{حرّك: C, (4, 6)} \\
            }
        }
    }
    & & \\ \\

    \theutterance \stepcounter{utterance}  
    & & & \multicolumn{2}{p{0.3\linewidth}}{
        \cellcolor[rgb]{0.9,0.9,0.9}{
            \makecell[{{p{\linewidth}}}]{
                \texttt{\tiny{[GM$|$GM]}}
                \texttt{C moved to (4, 6)} \\
            }
        }
    }
    & & \\ \\

    \theutterance \stepcounter{utterance}  
    & & \multicolumn{4}{p{0.6\linewidth}}{
        \cellcolor[rgb]{0.9,0.9,0.9}{
            \makecell[{{p{\linewidth}}}]{
                \texttt{\tiny{[GM$\rangle$P2]}}
                \texttt{تم تحريك 'T' إلى (3,3) بنجاح. تبدو شبكتك المحدّثة هكذا:} \\
\texttt{```} \\
\texttt{ 1234567} \\
\texttt{╔═╤═╤═══╗} \\
\texttt{║C│I│◌◌◌║ 1} \\
\texttt{╟─┴─┼─┬─╢ 2} \\
\texttt{║◌◌T│◌│◌║ 3} \\
\texttt{║◌◌◌├─┘◌║ 4} \\
\texttt{║◌┌─┘◌◌◌║ 5} \\
\texttt{╟─┤◌W◌◌H║ 6} \\
\texttt{║◌│◌◌◌◌◌║ 7} \\
\texttt{╚═╧═════╝} \\
\texttt{```} \\
\texttt{أنت تلعب حالياً الجولة 7 من أصل 20 كحد أقصى.} \\
\texttt{لقد جمعتما معاً 15 من أصل 12 عقوبة.} \\
\texttt{حرّك اللاعب الآخر كائناً على شبكته.} \\
\\ 
\\ 
\texttt{ما هو أمرك التالي؟} \\
            }
        }
    }
    & \\ \\

    \theutterance \stepcounter{utterance}  
    & & & \multicolumn{4}{p{0.6\linewidth}}{
        \cellcolor[rgb]{0.9,0.9,0.9}{
            \makecell[{{p{\linewidth}}}]{
                \texttt{\tiny{[GM$\langle$P2]}}
                \texttt{حرّك: T, (3, 1)} \\
            }
        }
    }
     \\ \\

    \theutterance \stepcounter{utterance}  
    & & & \multicolumn{2}{p{0.3\linewidth}}{
        \cellcolor[rgb]{0.9,0.9,0.9}{
            \makecell[{{p{\linewidth}}}]{
                \texttt{\tiny{[GM$|$GM]}}
                \texttt{حرّك: T, (3, 1)} \\
            }
        }
    }
    & & \\ \\

    \theutterance \stepcounter{utterance}  
    & & & \multicolumn{2}{p{0.3\linewidth}}{
        \cellcolor[rgb]{0.9,0.9,0.9}{
            \makecell[{{p{\linewidth}}}]{
                \texttt{\tiny{[GM$|$GM]}}
                \texttt{عقوبة: (3,1) ليست فارغة، بل تحتوي على 'I'.} \\
\\ 
\texttt{لقد جمعتما معاً 16 من أصل 12 عقوبة.} \\
\texttt{يرجى المحاولة مرة أخرى!} \\
            }
        }
    }
    & & \\ \\

    \theutterance \stepcounter{utterance}  
    & & \multicolumn{4}{p{0.6\linewidth}}{
        \cellcolor[rgb]{0.9,0.9,0.9}{
            \makecell[{{p{\linewidth}}}]{
                \texttt{\tiny{[GM$\rangle$P2]}}
                \texttt{عقوبة: (3,1) ليست فارغة، بل تحتوي على 'I'.} \\
\\ 
\texttt{لقد جمعتما معاً 16 من أصل 12 عقوبة.} \\
\texttt{يرجى المحاولة مرة أخرى!} \\
            }
        }
    }
    & \\ \\

    \theutterance \stepcounter{utterance}  
    & & & \multicolumn{4}{p{0.6\linewidth}}{
        \cellcolor[rgb]{0.9,0.9,0.9}{
            \makecell[{{p{\linewidth}}}]{
                \texttt{\tiny{[GM$\langle$P2]}}
                \texttt{حرّك: I, (2, 3)} \\
            }
        }
    }
     \\ \\

    \theutterance \stepcounter{utterance}  
    & & & \multicolumn{2}{p{0.3\linewidth}}{
        \cellcolor[rgb]{0.9,0.9,0.9}{
            \makecell[{{p{\linewidth}}}]{
                \texttt{\tiny{[GM$|$GM]}}
                \texttt{حرّك: I, (2, 3)} \\
            }
        }
    }
    & & \\ \\

    \theutterance \stepcounter{utterance}  
    & & & \multicolumn{2}{p{0.3\linewidth}}{
        \cellcolor[rgb]{0.9,0.9,0.9}{
            \makecell[{{p{\linewidth}}}]{
                \texttt{\tiny{[GM$|$GM]}}
                \texttt{I moved to (2, 3)} \\
            }
        }
    }
    & & \\ \\

    \theutterance \stepcounter{utterance}  
    & & \multicolumn{4}{p{0.6\linewidth}}{
        \cellcolor[rgb]{0.9,0.9,0.9}{
            \makecell[{{p{\linewidth}}}]{
                \texttt{\tiny{[P1$\langle$GM]}}
                \texttt{تم تحريك 'C' إلى (4,6) بنجاح. تبدو شبكتك المحدّثة هكذا:} \\
\texttt{```} \\
\texttt{ 1234567} \\
\texttt{╔═╤═╤═══╗} \\
\texttt{║W│◌│T◌◌║ 1} \\
\texttt{╟─┴─┼─┬─╢ 2} \\
\texttt{║IH◌│◌│◌║ 3} \\
\texttt{║◌◌◌├─┘◌║ 4} \\
\texttt{║◌┌─┘◌◌◌║ 5} \\
\texttt{╟─┤◌C◌◌◌║ 6} \\
\texttt{║◌│◌◌◌◌◌║ 7} \\
\texttt{╚═╧═════╝} \\
\texttt{```} \\
\texttt{أنت تلعب حالياً الجولة 7 من أصل 20 كحد أقصى.} \\
\texttt{لقد جمعتما معاً 16 من أصل 12 عقوبة.} \\
\texttt{حرّك اللاعب الآخر كائناً على شبكته.} \\
\\ 
\\ 
\texttt{ما هو أمرك التالي؟} \\
            }
        }
    }
    & \\ \\

    \theutterance \stepcounter{utterance}  
    & \multicolumn{4}{p{0.6\linewidth}}{
        \cellcolor[rgb]{0.9,0.9,0.9}{
            \makecell[{{p{\linewidth}}}]{
                \texttt{\tiny{[P1$\rangle$GM]}}
                \texttt{قل: لقد تجاوزنا عدد العقوبات المسموح به، يبدو أننا خسرنا اللعبة. أعتذر عن الأخطاء في تحديد الإحداثيات الفارغة.} \\
            }
        }
    }
    & & \\ \\

    \theutterance \stepcounter{utterance}  
    & & & \multicolumn{2}{p{0.3\linewidth}}{
        \cellcolor[rgb]{0.9,0.9,0.9}{
            \makecell[{{p{\linewidth}}}]{
                \texttt{\tiny{[GM$|$GM]}}
                \texttt{قل: لقد تجاوزنا عدد العقوبات المسموح به، يبدو أننا خسرنا اللعبة. أعتذر عن الأخطاء في تحديد الإحداثيات الفارغة.} \\
            }
        }
    }
    & & \\ \\

    \theutterance \stepcounter{utterance}  
    & & \multicolumn{4}{p{0.6\linewidth}}{
        \cellcolor[rgb]{0.9,0.9,0.9}{
            \makecell[{{p{\linewidth}}}]{
                \texttt{\tiny{[GM$\rangle$P2]}}
                \texttt{تم تحريك 'I' إلى (2,3) بنجاح. تبدو شبكتك المحدّثة هكذا:} \\
\texttt{```} \\
\texttt{ 1234567} \\
\texttt{╔═╤═╤═══╗} \\
\texttt{║C│◌│◌◌◌║ 1} \\
\texttt{╟─┴─┼─┬─╢ 2} \\
\texttt{║◌IT│◌│◌║ 3} \\
\texttt{║◌◌◌├─┘◌║ 4} \\
\texttt{║◌┌─┘◌◌◌║ 5} \\
\texttt{╟─┤◌W◌◌H║ 6} \\
\texttt{║◌│◌◌◌◌◌║ 7} \\
\texttt{╚═╧═════╝} \\
\texttt{```} \\
\texttt{أنت تلعب حالياً الجولة 8 من أصل 20 كحد أقصى.} \\
\texttt{لقد جمعتما معاً 16 من أصل 12 عقوبة.} \\
\texttt{أرسل اللاعب الآخر هذه الرسالة:} \\
\\ 
\\ 
\texttt{"لقد تجاوزنا عدد العقوبات المسموح به، يبدو أننا خسرنا اللعبة. أعتذر عن الأخطاء في تحديد الإحداثيات الفارغة."} \\
\\ 
\\ 
\texttt{ما هو أمرك التالي؟} \\
            }
        }
    }
    & \\ \\

    \theutterance \stepcounter{utterance}  
    & & & \multicolumn{4}{p{0.6\linewidth}}{
        \cellcolor[rgb]{0.9,0.9,0.9}{
            \makecell[{{p{\linewidth}}}]{
                \texttt{\tiny{[GM$\langle$P2]}}
                \texttt{قل: لا بأس، لقد كانت تجربة تعليمية جيدة. أعتذر أيضاً عن الأخطاء في قراءة الشبكة.} \\
            }
        }
    }
     \\ \\

    \theutterance \stepcounter{utterance}  
    & & & \multicolumn{2}{p{0.3\linewidth}}{
        \cellcolor[rgb]{0.9,0.9,0.9}{
            \makecell[{{p{\linewidth}}}]{
                \texttt{\tiny{[GM$|$GM]}}
                \texttt{قل: لا بأس، لقد كانت تجربة تعليمية جيدة. أعتذر أيضاً عن الأخطاء في قراءة الشبكة.} \\
            }
        }
    }
    & & \\ \\

    \theutterance \stepcounter{utterance}  
    & & \multicolumn{4}{p{0.6\linewidth}}{
        \cellcolor[rgb]{0.9,0.9,0.9}{
            \makecell[{{p{\linewidth}}}]{
                \texttt{\tiny{[P1$\langle$GM]}}
                \texttt{تم تمرير رسالتك إلى اللاعب الآخر.} \\
\\ 
\\ 
\texttt{أنت تلعب حالياً الجولة 8 من أصل 20 كحد أقصى.} \\
\texttt{لقد جمعتما معاً 16 من أصل 12 عقوبة.} \\
\texttt{أرسل اللاعب الآخر هذه الرسالة:} \\
\\ 
\\ 
\texttt{"لا بأس، لقد كانت تجربة تعليمية جيدة. أعتذر أيضاً عن الأخطاء في قراءة الشبكة."} \\
\\ 
\\ 
\texttt{ما هو أمرك التالي؟} \\
            }
        }
    }
    & \\ \\

    \theutterance \stepcounter{utterance}  
    & \multicolumn{4}{p{0.6\linewidth}}{
        \cellcolor[rgb]{0.9,0.9,0.9}{
            \makecell[{{p{\linewidth}}}]{
                \texttt{\tiny{[P1$\rangle$GM]}}
                \texttt{قل: لا بأس، لنحاول الاستمرار إذا كان بإمكاننا ذلك، رغم أننا تجاوزنا الحد. سأحاول تحريك الكائنات المتبقية بدقة أكبر.} \\
            }
        }
    }
    & & \\ \\

    \theutterance \stepcounter{utterance}  
    & & & \multicolumn{2}{p{0.3\linewidth}}{
        \cellcolor[rgb]{0.9,0.9,0.9}{
            \makecell[{{p{\linewidth}}}]{
                \texttt{\tiny{[GM$|$GM]}}
                \texttt{قل: لا بأس، لنحاول الاستمرار إذا كان بإمكاننا ذلك، رغم أننا تجاوزنا الحد. سأحاول تحريك الكائنات المتبقية بدقة أكبر.} \\
            }
        }
    }
    & & \\ \\

    \theutterance \stepcounter{utterance}  
    & & \multicolumn{4}{p{0.6\linewidth}}{
        \cellcolor[rgb]{0.9,0.9,0.9}{
            \makecell[{{p{\linewidth}}}]{
                \texttt{\tiny{[GM$\rangle$P2]}}
                \texttt{تم تمرير رسالتك إلى اللاعب الآخر.} \\
\\ 
\\ 
\texttt{أنت تلعب حالياً الجولة 9 من أصل 20 كحد أقصى.} \\
\texttt{لقد جمعتما معاً 16 من أصل 12 عقوبة.} \\
\texttt{أرسل اللاعب الآخر هذه الرسالة:} \\
\\ 
\\ 
\texttt{"لا بأس، لنحاول الاستمرار إذا كان بإمكاننا ذلك، رغم أننا تجاوزنا الحد. سأحاول تحريك الكائنات المتبقية بدقة أكبر."} \\
\\ 
\\ 
\texttt{ما هو أمرك التالي؟} \\
            }
        }
    }
    & \\ \\

    \theutterance \stepcounter{utterance}  
    & & & \multicolumn{4}{p{0.6\linewidth}}{
        \cellcolor[rgb]{0.9,0.9,0.9}{
            \makecell[{{p{\linewidth}}}]{
                \texttt{\tiny{[GM$\langle$P2]}}
                \texttt{قل: حسناً، لنحاول المضي قدماً بكل حذر. سأركز على تحريك الكائنات بدقة شديدة لتجنب أي عقوبات إضافية.} \\
            }
        }
    }
     \\ \\

    \theutterance \stepcounter{utterance}  
    & & & \multicolumn{2}{p{0.3\linewidth}}{
        \cellcolor[rgb]{0.9,0.9,0.9}{
            \makecell[{{p{\linewidth}}}]{
                \texttt{\tiny{[GM$|$GM]}}
                \texttt{قل: حسناً، لنحاول المضي قدماً بكل حذر. سأركز على تحريك الكائنات بدقة شديدة لتجنب أي عقوبات إضافية.} \\
            }
        }
    }
    & & \\ \\

    \theutterance \stepcounter{utterance}  
    & & \multicolumn{4}{p{0.6\linewidth}}{
        \cellcolor[rgb]{0.9,0.9,0.9}{
            \makecell[{{p{\linewidth}}}]{
                \texttt{\tiny{[P1$\langle$GM]}}
                \texttt{تم تمرير رسالتك إلى اللاعب الآخر.} \\
\\ 
\\ 
\texttt{أنت تلعب حالياً الجولة 9 من أصل 20 كحد أقصى.} \\
\texttt{لقد جمعتما معاً 16 من أصل 12 عقوبة.} \\
\texttt{أرسل اللاعب الآخر هذه الرسالة:} \\
\\ 
\\ 
\texttt{"حسناً، لنحاول المضي قدماً بكل حذر. سأركز على تحريك الكائنات بدقة شديدة لتجنب أي عقوبات إضافية."} \\
\\ 
\\ 
\texttt{ما هو أمرك التالي؟} \\
            }
        }
    }
    & \\ \\

    \theutterance \stepcounter{utterance}  
    & \multicolumn{4}{p{0.6\linewidth}}{
        \cellcolor[rgb]{0.9,0.9,0.9}{
            \makecell[{{p{\linewidth}}}]{
                \texttt{\tiny{[P1$\rangle$GM]}}
                \texttt{حرّك: C, (5, 6)} \\
            }
        }
    }
    & & \\ \\

    \theutterance \stepcounter{utterance}  
    & & & \multicolumn{2}{p{0.3\linewidth}}{
        \cellcolor[rgb]{0.9,0.9,0.9}{
            \makecell[{{p{\linewidth}}}]{
                \texttt{\tiny{[GM$|$GM]}}
                \texttt{حرّك: C, (5, 6)} \\
            }
        }
    }
    & & \\ \\

    \theutterance \stepcounter{utterance}  
    & & & \multicolumn{2}{p{0.3\linewidth}}{
        \cellcolor[rgb]{0.9,0.9,0.9}{
            \makecell[{{p{\linewidth}}}]{
                \texttt{\tiny{[GM$|$GM]}}
                \texttt{C moved to (5, 6)} \\
            }
        }
    }
    & & \\ \\

    \theutterance \stepcounter{utterance}  
    & & \multicolumn{4}{p{0.6\linewidth}}{
        \cellcolor[rgb]{0.9,0.9,0.9}{
            \makecell[{{p{\linewidth}}}]{
                \texttt{\tiny{[GM$\rangle$P2]}}
                \texttt{تم تمرير رسالتك إلى اللاعب الآخر.} \\
\\ 
\\ 
\texttt{أنت تلعب حالياً الجولة 10 من أصل 20 كحد أقصى.} \\
\texttt{لقد جمعتما معاً 16 من أصل 12 عقوبة.} \\
\texttt{حرّك اللاعب الآخر كائناً على شبكته.} \\
\\ 
\\ 
\texttt{ما هو أمرك التالي؟} \\
            }
        }
    }
    & \\ \\

    \theutterance \stepcounter{utterance}  
    & & & \multicolumn{4}{p{0.6\linewidth}}{
        \cellcolor[rgb]{0.9,0.9,0.9}{
            \makecell[{{p{\linewidth}}}]{
                \texttt{\tiny{[GM$\langle$P2]}}
                \texttt{حرّك: T, (3, 1)} \\
            }
        }
    }
     \\ \\

    \theutterance \stepcounter{utterance}  
    & & & \multicolumn{2}{p{0.3\linewidth}}{
        \cellcolor[rgb]{0.9,0.9,0.9}{
            \makecell[{{p{\linewidth}}}]{
                \texttt{\tiny{[GM$|$GM]}}
                \texttt{حرّك: T, (3, 1)} \\
            }
        }
    }
    & & \\ \\

    \theutterance \stepcounter{utterance}  
    & & & \multicolumn{2}{p{0.3\linewidth}}{
        \cellcolor[rgb]{0.9,0.9,0.9}{
            \makecell[{{p{\linewidth}}}]{
                \texttt{\tiny{[GM$|$GM]}}
                \texttt{T moved to (3, 1)} \\
            }
        }
    }
    & & \\ \\

    \theutterance \stepcounter{utterance}  
    & & \multicolumn{4}{p{0.6\linewidth}}{
        \cellcolor[rgb]{0.9,0.9,0.9}{
            \makecell[{{p{\linewidth}}}]{
                \texttt{\tiny{[P1$\langle$GM]}}
                \texttt{تم تحريك 'C' إلى (5,6) بنجاح. تبدو شبكتك المحدّثة هكذا:} \\
\texttt{```} \\
\texttt{ 1234567} \\
\texttt{╔═╤═╤═══╗} \\
\texttt{║W│◌│T◌◌║ 1} \\
\texttt{╟─┴─┼─┬─╢ 2} \\
\texttt{║IH◌│◌│◌║ 3} \\
\texttt{║◌◌◌├─┘◌║ 4} \\
\texttt{║◌┌─┘◌◌◌║ 5} \\
\texttt{╟─┤◌◌C◌◌║ 6} \\
\texttt{║◌│◌◌◌◌◌║ 7} \\
\texttt{╚═╧═════╝} \\
\texttt{```} \\
\texttt{أنت تلعب حالياً الجولة 10 من أصل 20 كحد أقصى.} \\
\texttt{لقد جمعتما معاً 16 من أصل 12 عقوبة.} \\
\texttt{حرّك اللاعب الآخر كائناً على شبكته.} \\
\\ 
\\ 
\texttt{ما هو أمرك التالي؟} \\
            }
        }
    }
    & \\ \\

    \theutterance \stepcounter{utterance}  
    & \multicolumn{4}{p{0.6\linewidth}}{
        \cellcolor[rgb]{0.9,0.9,0.9}{
            \makecell[{{p{\linewidth}}}]{
                \texttt{\tiny{[P1$\rangle$GM]}}
                \texttt{حرّك: T, (6, 1)} \\
            }
        }
    }
    & & \\ \\

    \theutterance \stepcounter{utterance}  
    & & & \multicolumn{2}{p{0.3\linewidth}}{
        \cellcolor[rgb]{0.9,0.9,0.9}{
            \makecell[{{p{\linewidth}}}]{
                \texttt{\tiny{[GM$|$GM]}}
                \texttt{حرّك: T, (6, 1)} \\
            }
        }
    }
    & & \\ \\

    \theutterance \stepcounter{utterance}  
    & & & \multicolumn{2}{p{0.3\linewidth}}{
        \cellcolor[rgb]{0.9,0.9,0.9}{
            \makecell[{{p{\linewidth}}}]{
                \texttt{\tiny{[GM$|$GM]}}
                \texttt{T moved to (6, 1)} \\
            }
        }
    }
    & & \\ \\

    \theutterance \stepcounter{utterance}  
    & & \multicolumn{4}{p{0.6\linewidth}}{
        \cellcolor[rgb]{0.9,0.9,0.9}{
            \makecell[{{p{\linewidth}}}]{
                \texttt{\tiny{[GM$\rangle$P2]}}
                \texttt{تم تحريك 'T' إلى (3,1) بنجاح. تبدو شبكتك المحدّثة هكذا:} \\
\texttt{```} \\
\texttt{ 1234567} \\
\texttt{╔═╤═╤═══╗} \\
\texttt{║C│T│◌◌◌║ 1} \\
\texttt{╟─┴─┼─┬─╢ 2} \\
\texttt{║◌I◌│◌│◌║ 3} \\
\texttt{║◌◌◌├─┘◌║ 4} \\
\texttt{║◌┌─┘◌◌◌║ 5} \\
\texttt{╟─┤◌W◌◌H║ 6} \\
\texttt{║◌│◌◌◌◌◌║ 7} \\
\texttt{╚═╧═════╝} \\
\texttt{```} \\
\texttt{أنت تلعب حالياً الجولة 11 من أصل 20 كحد أقصى.} \\
\texttt{لقد جمعتما معاً 16 من أصل 12 عقوبة.} \\
\texttt{حرّك اللاعب الآخر كائناً على شبكته.} \\
\\ 
\\ 
\texttt{ما هو أمرك التالي؟} \\
            }
        }
    }
    & \\ \\

    \theutterance \stepcounter{utterance}  
    & & & \multicolumn{4}{p{0.6\linewidth}}{
        \cellcolor[rgb]{0.9,0.9,0.9}{
            \makecell[{{p{\linewidth}}}]{
                \texttt{\tiny{[GM$\langle$P2]}}
                \texttt{حرّك: I, (2, 1)} \\
            }
        }
    }
     \\ \\

    \theutterance \stepcounter{utterance}  
    & & & \multicolumn{2}{p{0.3\linewidth}}{
        \cellcolor[rgb]{0.9,0.9,0.9}{
            \makecell[{{p{\linewidth}}}]{
                \texttt{\tiny{[GM$|$GM]}}
                \texttt{حرّك: I, (2, 1)} \\
            }
        }
    }
    & & \\ \\

    \theutterance \stepcounter{utterance}  
    & & & \multicolumn{2}{p{0.3\linewidth}}{
        \cellcolor[rgb]{0.9,0.9,0.9}{
            \makecell[{{p{\linewidth}}}]{
                \texttt{\tiny{[GM$|$GM]}}
                \texttt{عقوبة: (2,1) ليست فارغة، بل تحتوي على '│'.} \\
\\ 
\texttt{لقد جمعتما معاً 17 من أصل 12 عقوبة.} \\
\texttt{يرجى المحاولة مرة أخرى!} \\
            }
        }
    }
    & & \\ \\

    \theutterance \stepcounter{utterance}  
    & & \multicolumn{4}{p{0.6\linewidth}}{
        \cellcolor[rgb]{0.9,0.9,0.9}{
            \makecell[{{p{\linewidth}}}]{
                \texttt{\tiny{[GM$\rangle$P2]}}
                \texttt{عقوبة: (2,1) ليست فارغة، بل تحتوي على '│'.} \\
\\ 
\texttt{لقد جمعتما معاً 17 من أصل 12 عقوبة.} \\
\texttt{يرجى المحاولة مرة أخرى!} \\
            }
        }
    }
    & \\ \\

    \theutterance \stepcounter{utterance}  
    & & & \multicolumn{4}{p{0.6\linewidth}}{
        \cellcolor[rgb]{0.9,0.9,0.9}{
            \makecell[{{p{\linewidth}}}]{
                \texttt{\tiny{[GM$\langle$P2]}}
                \texttt{حرّك: I, (2, 3)} \\
            }
        }
    }
     \\ \\

    \theutterance \stepcounter{utterance}  
    & & & \multicolumn{2}{p{0.3\linewidth}}{
        \cellcolor[rgb]{0.9,0.9,0.9}{
            \makecell[{{p{\linewidth}}}]{
                \texttt{\tiny{[GM$|$GM]}}
                \texttt{حرّك: I, (2, 3)} \\
            }
        }
    }
    & & \\ \\

    \theutterance \stepcounter{utterance}  
    & & & \multicolumn{2}{p{0.3\linewidth}}{
        \cellcolor[rgb]{0.9,0.9,0.9}{
            \makecell[{{p{\linewidth}}}]{
                \texttt{\tiny{[GM$|$GM]}}
                \texttt{عقوبة: (2,3) ليست فارغة، بل تحتوي على 'I'.} \\
\\ 
\texttt{لقد جمعتما معاً 18 من أصل 12 عقوبة.} \\
\texttt{يرجى المحاولة مرة أخرى!} \\
            }
        }
    }
    & & \\ \\

    \theutterance \stepcounter{utterance}  
    & & \multicolumn{4}{p{0.6\linewidth}}{
        \cellcolor[rgb]{0.9,0.9,0.9}{
            \makecell[{{p{\linewidth}}}]{
                \texttt{\tiny{[GM$\rangle$P2]}}
                \texttt{عقوبة: (2,3) ليست فارغة، بل تحتوي على 'I'.} \\
\\ 
\texttt{لقد جمعتما معاً 18 من أصل 12 عقوبة.} \\
\texttt{يرجى المحاولة مرة أخرى!} \\
            }
        }
    }
    & \\ \\

    \theutterance \stepcounter{utterance}  
    & & & \multicolumn{4}{p{0.6\linewidth}}{
        \cellcolor[rgb]{0.9,0.9,0.9}{
            \makecell[{{p{\linewidth}}}]{
                \texttt{\tiny{[GM$\langle$P2]}}
                \texttt{حرّك: I, (1, 3)} \\
            }
        }
    }
     \\ \\

    \theutterance \stepcounter{utterance}  
    & & & \multicolumn{2}{p{0.3\linewidth}}{
        \cellcolor[rgb]{0.9,0.9,0.9}{
            \makecell[{{p{\linewidth}}}]{
                \texttt{\tiny{[GM$|$GM]}}
                \texttt{حرّك: I, (1, 3)} \\
            }
        }
    }
    & & \\ \\

    \theutterance \stepcounter{utterance}  
    & & & \multicolumn{2}{p{0.3\linewidth}}{
        \cellcolor[rgb]{0.9,0.9,0.9}{
            \makecell[{{p{\linewidth}}}]{
                \texttt{\tiny{[GM$|$GM]}}
                \texttt{I moved to (1, 3)} \\
            }
        }
    }
    & & \\ \\

    \theutterance \stepcounter{utterance}  
    & & \multicolumn{4}{p{0.6\linewidth}}{
        \cellcolor[rgb]{0.9,0.9,0.9}{
            \makecell[{{p{\linewidth}}}]{
                \texttt{\tiny{[P1$\langle$GM]}}
                \texttt{تم تحريك 'T' إلى (6,1) بنجاح. تبدو شبكتك المحدّثة هكذا:} \\
\texttt{```} \\
\texttt{ 1234567} \\
\texttt{╔═╤═╤═══╗} \\
\texttt{║W│◌│◌T◌║ 1} \\
\texttt{╟─┴─┼─┬─╢ 2} \\
\texttt{║IH◌│◌│◌║ 3} \\
\texttt{║◌◌◌├─┘◌║ 4} \\
\texttt{║◌┌─┘◌◌◌║ 5} \\
\texttt{╟─┤◌◌C◌◌║ 6} \\
\texttt{║◌│◌◌◌◌◌║ 7} \\
\texttt{╚═╧═════╝} \\
\texttt{```} \\
\texttt{أنت تلعب حالياً الجولة 11 من أصل 20 كحد أقصى.} \\
\texttt{لقد جمعتما معاً 18 من أصل 12 عقوبة.} \\
\texttt{حرّك اللاعب الآخر كائناً على شبكته.} \\
\\ 
\\ 
\texttt{ما هو أمرك التالي؟} \\
            }
        }
    }
    & \\ \\

    \theutterance \stepcounter{utterance}  
    & \multicolumn{4}{p{0.6\linewidth}}{
        \cellcolor[rgb]{0.9,0.9,0.9}{
            \makecell[{{p{\linewidth}}}]{
                \texttt{\tiny{[P1$\rangle$GM]}}
                \texttt{حرّك: H, (2, 5)} \\
            }
        }
    }
    & & \\ \\

    \theutterance \stepcounter{utterance}  
    & & & \multicolumn{2}{p{0.3\linewidth}}{
        \cellcolor[rgb]{0.9,0.9,0.9}{
            \makecell[{{p{\linewidth}}}]{
                \texttt{\tiny{[GM$|$GM]}}
                \texttt{حرّك: H, (2, 5)} \\
            }
        }
    }
    & & \\ \\

    \theutterance \stepcounter{utterance}  
    & & & \multicolumn{2}{p{0.3\linewidth}}{
        \cellcolor[rgb]{0.9,0.9,0.9}{
            \makecell[{{p{\linewidth}}}]{
                \texttt{\tiny{[GM$|$GM]}}
                \texttt{عقوبة: (2,5) ليست فارغة، بل تحتوي على '┌'.} \\
\\ 
\texttt{لقد جمعتما معاً 19 من أصل 12 عقوبة.} \\
\texttt{يرجى المحاولة مرة أخرى!} \\
            }
        }
    }
    & & \\ \\

    \theutterance \stepcounter{utterance}  
    & & \multicolumn{4}{p{0.6\linewidth}}{
        \cellcolor[rgb]{0.9,0.9,0.9}{
            \makecell[{{p{\linewidth}}}]{
                \texttt{\tiny{[P1$\langle$GM]}}
                \texttt{عقوبة: (2,5) ليست فارغة، بل تحتوي على '┌'.} \\
\\ 
\texttt{لقد جمعتما معاً 19 من أصل 12 عقوبة.} \\
\texttt{يرجى المحاولة مرة أخرى!} \\
            }
        }
    }
    & \\ \\

    \theutterance \stepcounter{utterance}  
    & \multicolumn{4}{p{0.6\linewidth}}{
        \cellcolor[rgb]{0.9,0.9,0.9}{
            \makecell[{{p{\linewidth}}}]{
                \texttt{\tiny{[P1$\rangle$GM]}}
                \texttt{حرّك: H, (3, 5)} \\
            }
        }
    }
    & & \\ \\

    \theutterance \stepcounter{utterance}  
    & & & \multicolumn{2}{p{0.3\linewidth}}{
        \cellcolor[rgb]{0.9,0.9,0.9}{
            \makecell[{{p{\linewidth}}}]{
                \texttt{\tiny{[GM$|$GM]}}
                \texttt{حرّك: H, (3, 5)} \\
            }
        }
    }
    & & \\ \\

    \theutterance \stepcounter{utterance}  
    & & & \multicolumn{2}{p{0.3\linewidth}}{
        \cellcolor[rgb]{0.9,0.9,0.9}{
            \makecell[{{p{\linewidth}}}]{
                \texttt{\tiny{[GM$|$GM]}}
                \texttt{عقوبة: (3,5) ليست فارغة، بل تحتوي على '─'.} \\
\\ 
\texttt{لقد جمعتما معاً 20 من أصل 12 عقوبة.} \\
\texttt{يرجى المحاولة مرة أخرى!} \\
            }
        }
    }
    & & \\ \\

    \theutterance \stepcounter{utterance}  
    & & \multicolumn{4}{p{0.6\linewidth}}{
        \cellcolor[rgb]{0.9,0.9,0.9}{
            \makecell[{{p{\linewidth}}}]{
                \texttt{\tiny{[P1$\langle$GM]}}
                \texttt{عقوبة: (3,5) ليست فارغة، بل تحتوي على '─'.} \\
\\ 
\texttt{لقد جمعتما معاً 20 من أصل 12 عقوبة.} \\
\texttt{يرجى المحاولة مرة أخرى!} \\
            }
        }
    }
    & \\ \\

    \theutterance \stepcounter{utterance}  
    & \multicolumn{4}{p{0.6\linewidth}}{
        \cellcolor[rgb]{0.9,0.9,0.9}{
            \makecell[{{p{\linewidth}}}]{
                \texttt{\tiny{[P1$\rangle$GM]}}
                \texttt{حرّك: H, (2, 4)} \\
            }
        }
    }
    & & \\ \\

    \theutterance \stepcounter{utterance}  
    & & & \multicolumn{2}{p{0.3\linewidth}}{
        \cellcolor[rgb]{0.9,0.9,0.9}{
            \makecell[{{p{\linewidth}}}]{
                \texttt{\tiny{[GM$|$GM]}}
                \texttt{حرّك: H, (2, 4)} \\
            }
        }
    }
    & & \\ \\

    \theutterance \stepcounter{utterance}  
    & & & \multicolumn{2}{p{0.3\linewidth}}{
        \cellcolor[rgb]{0.9,0.9,0.9}{
            \makecell[{{p{\linewidth}}}]{
                \texttt{\tiny{[GM$|$GM]}}
                \texttt{H moved to (2, 4)} \\
            }
        }
    }
    & & \\ \\

    \theutterance \stepcounter{utterance}  
    & & \multicolumn{4}{p{0.6\linewidth}}{
        \cellcolor[rgb]{0.9,0.9,0.9}{
            \makecell[{{p{\linewidth}}}]{
                \texttt{\tiny{[GM$\rangle$P2]}}
                \texttt{تم تحريك 'I' إلى (1,3) بنجاح. تبدو شبكتك المحدّثة هكذا:} \\
\texttt{```} \\
\texttt{ 1234567} \\
\texttt{╔═╤═╤═══╗} \\
\texttt{║C│T│◌◌◌║ 1} \\
\texttt{╟─┴─┼─┬─╢ 2} \\
\texttt{║I◌◌│◌│◌║ 3} \\
\texttt{║◌◌◌├─┘◌║ 4} \\
\texttt{║◌┌─┘◌◌◌║ 5} \\
\texttt{╟─┤◌W◌◌H║ 6} \\
\texttt{║◌│◌◌◌◌◌║ 7} \\
\texttt{╚═╧═════╝} \\
\texttt{```} \\
\texttt{أنت تلعب حالياً الجولة 12 من أصل 20 كحد أقصى.} \\
\texttt{لقد جمعتما معاً 20 من أصل 12 عقوبة.} \\
\texttt{حرّك اللاعب الآخر كائناً على شبكته.} \\
\\ 
\\ 
\texttt{ما هو أمرك التالي؟} \\
            }
        }
    }
    & \\ \\

    \theutterance \stepcounter{utterance}  
    & & & \multicolumn{4}{p{0.6\linewidth}}{
        \cellcolor[rgb]{0.9,0.9,0.9}{
            \makecell[{{p{\linewidth}}}]{
                \texttt{\tiny{[GM$\langle$P2]}}
                \texttt{حرّك: I, (2, 1)} \\
\\ 
\texttt{(ملاحظة: سأحاول تحريك I إلى (2,1) مرة أخرى، ولكن بما أن (2,1) تحتوي على '│'، سأختار موضعاً فارغاً تماماً. بالنظر للشبكة، (2,3) تبدو فارغة `◌`).} \\
\\ 
\texttt{حرّك: I, (2, 3)} \\
            }
        }
    }
     \\ \\

    \theutterance \stepcounter{utterance}  
    & & & \multicolumn{2}{p{0.3\linewidth}}{
        \cellcolor[rgb]{0.9,0.9,0.9}{
            \makecell[{{p{\linewidth}}}]{
                \texttt{\tiny{[GM$|$GM]}}
                \texttt{حرّك: I, (2, 1)} \\
\\ 
\texttt{(ملاحظة: سأحاول تحريك I إلى (2,1) مرة أخرى، ولكن بما أن (2,1) تحتوي على '│'، سأختار موضعاً فارغاً تماماً. بالنظر للشبكة، (2,3) تبدو فارغة `◌`).} \\
\\ 
\texttt{حرّك: I, (2, 3)} \\
            }
        }
    }
    & & \\ \\

    \theutterance \stepcounter{utterance}  
    & & & \multicolumn{2}{p{0.3\linewidth}}{
        \cellcolor[rgb]{0.9,0.9,0.9}{
            \makecell[{{p{\linewidth}}}]{
                \texttt{\tiny{[GM$|$GM]}}
                \texttt{Invalid format: head is not empty: 'حرّك: I, (2, 1)} \\
\\ 
\texttt{(ملاحظة: سأحاول تحريك I إلى (2,1) مرة أخرى، ولكن بما أن (2,1) تحتوي على '│'، سأختار موضعاً فارغاً تماماً. بالنظر للشبكة، (2,3) تبدو فارغة ◌).} \\
\\ 
\texttt{'} \\
            }
        }
    }
    & & \\ \\

    \theutterance \stepcounter{utterance}  
    & & & \multicolumn{2}{p{0.3\linewidth}}{
        \cellcolor[rgb]{0.9,0.9,0.9}{
            \makecell[{{p{\linewidth}}}]{
                \texttt{\tiny{[GM$|$GM]}}
                \texttt{Maximum number of penalties exceeded} \\
            }
        }
    }
    & & \\ \\

    \theutterance \stepcounter{utterance}  
    & & & \multicolumn{2}{p{0.3\linewidth}}{
        \cellcolor[rgb]{0.9,0.9,0.9}{
            \makecell[{{p{\linewidth}}}]{
                \texttt{\tiny{[GM$|$GM]}}
                \texttt{Initial states:} \\
\texttt{```} \\
\texttt{Player 1:} \\
\texttt{ 1234567} \\
\texttt{╔═╤═╤═══╗} \\
\texttt{║◌│◌│◌◌◌║ 1} \\
\texttt{╟─┴─┼─┬─╢ 2} \\
\texttt{║◌◌◌│◌│T║ 3} \\
\texttt{║◌H◌├─┘◌║ 4} \\
\texttt{║◌┌─┘◌◌W║ 5} \\
\texttt{╟─┤◌◌◌◌◌║ 6} \\
\texttt{║I│◌C◌◌◌║ 7} \\
\texttt{╚═╧═════╝} \\
\\ 
\\ 
\texttt{Player 2:} \\
\texttt{ 1234567} \\
\texttt{╔═╤═╤═══╗} \\
\texttt{║◌│◌│◌◌◌║ 1} \\
\texttt{╟─┴─┼─┬─╢ 2} \\
\texttt{║◌◌◌│◌│◌║ 3} \\
\texttt{║◌◌C├─┘◌║ 4} \\
\texttt{║◌┌─┘◌T◌║ 5} \\
\texttt{╟─┤◌W◌◌H║ 6} \\
\texttt{║◌│◌I◌◌◌║ 7} \\
\texttt{╚═╧═════╝} \\
\\ 
\texttt{```} \\
            }
        }
    }
    & & \\ \\

    \theutterance \stepcounter{utterance}  
    & & & \multicolumn{2}{p{0.3\linewidth}}{
        \cellcolor[rgb]{0.9,0.9,0.9}{
            \makecell[{{p{\linewidth}}}]{
                \texttt{\tiny{[GM$|$GM]}}
                \texttt{End states:} \\
\texttt{```} \\
\texttt{Player 1:} \\
\texttt{ 1234567} \\
\texttt{╔═╤═╤═══╗} \\
\texttt{║W│◌│◌T◌║ 1} \\
\texttt{╟─┴─┼─┬─╢ 2} \\
\texttt{║I◌◌│◌│◌║ 3} \\
\texttt{║◌H◌├─┘◌║ 4} \\
\texttt{║◌┌─┘◌◌◌║ 5} \\
\texttt{╟─┤◌◌C◌◌║ 6} \\
\texttt{║◌│◌◌◌◌◌║ 7} \\
\texttt{╚═╧═════╝} \\
\\ 
\\ 
\texttt{Player 2:} \\
\texttt{ 1234567} \\
\texttt{╔═╤═╤═══╗} \\
\texttt{║C│T│◌◌◌║ 1} \\
\texttt{╟─┴─┼─┬─╢ 2} \\
\texttt{║I◌◌│◌│◌║ 3} \\
\texttt{║◌◌◌├─┘◌║ 4} \\
\texttt{║◌┌─┘◌◌◌║ 5} \\
\texttt{╟─┤◌W◌◌H║ 6} \\
\texttt{║◌│◌◌◌◌◌║ 7} \\
\texttt{╚═╧═════╝} \\
\\ 
\texttt{```} \\
            }
        }
    }
    & & \\ \\

    \theutterance \stepcounter{utterance}  
    & & & \multicolumn{2}{p{0.3\linewidth}}{
        \cellcolor[rgb]{0.9,0.9,0.9}{
            \makecell[{{p{\linewidth}}}]{
                \texttt{\tiny{[GM$|$GM]}}
                \texttt{* success: False} \\
\texttt{* lose: False} \\
\texttt{* aborted: True} \\
\texttt{{-}{-}{-}{-}{-}{-}{-}} \\
\texttt{* Shifts: 13.00} \\
\texttt{* Max Shifts: 8.00} \\
\texttt{* Min Shifts: 4.00} \\
\texttt{* End Distance Sum: 20.62} \\
\texttt{* Init Distance Sum: 16.95} \\
\texttt{* Expected Distance Sum: 20.95} \\
\texttt{* Penalties: 21.00} \\
\texttt{* Max Penalties: 12.00} \\
\texttt{* Rounds: 11.00} \\
\texttt{* Max Rounds: 20.00} \\
\texttt{* Object Count: 5.00} \\
            }
        }
    }
    & & \\ \\

\end{supertabular}
}

\end{document}
